% !TeX root = ../../Book2.tex
% 纲要标题：## 第15章　嵌入、邻域与光滑族的局部平凡性
% 纲要位置：几何学计划纲要.md，第 1572 行。
% 正文状态：初稿；尚未正式审核。

% 纲要内容开始
% 
% > 本章目标与证明范围：将嵌入、管状邻域、同位延拓与 Ehresmann 按所需局部模型连接。弱嵌入维数界与强 Whitney 界分别安排证明，不能用一次泛型投影跨过消去双点的步骤。
% 
% 纲要内容结束

\chapter{嵌入、邻域与光滑族的局部平凡性}
\label{g2:ch:15}
\BookChapterNumber{g2:ch:15}

\begin{chapterabstract}
	嵌入把抽象流形放进 Euclid 空间，管状邻域把子流形附近的几何化为法丛，领状邻域则把边界附近化为乘积。本章先由坐标和泛型投影建立弱 Whitney 界，再补充双点消去与低维构造，得到强嵌入界。随后用逆函数定理、变半径控制和流构造邻域模型，并把子流形的同位延拓为环境运动。最后证明固有满淹没局部平凡：水平提升给出纤维间的搬运，固有性保证提升存在到所需时刻。边界条件和复环面族说明，局部平凡性的结论取决于所处的光滑或复解析范畴。
\end{chapterabstract}

\begin{learninggoals}
	\begin{itemize}
		\item 从有限坐标与光滑截断构造初始嵌入，计算投影坏方向的维数，并解释非紧源为何还需固有性控制。
		\item 辨认一般位置中的横截双点，说明 Whitney 盘、标架及局部移动如何参与双点消去，并分清低维情形。
		\item 用法向局部指数和变半径缩小构造管状邻域，理解法丛正规化在选择无关性中的作用。
		\item 从向内横截向量场构造领圈，利用有符号的领圈参数建立双化与边界粘合。
		\item 把同位的速度延拓为时间依赖向量场，用共同紧支集获得环境同位，并识别相对条件和非紧限制。
		\item 由水平提升证明 Ehresmann 定理，核对边界推广条件，并区分光滑局部平凡与复解析局部平凡。
	\end{itemize}
\end{learninggoals}

一条光滑曲线在每个点附近都像直线，一张曲面在每个点附近都像平面。这种局部坐标并没有立即给出整条曲线或整张曲面的统一表达。例如，沿一个圆周逐点选择短法线段，必须保证来自不同基点的线段不会相撞，才能用“圆周上的位置加法向位移”描述整个邻域；若圆周随时间运动，还要把不同位置上的速度组织成一个环境向量场。本章研究的正是这些从逐点局部模型到整体构造的步骤。

\ProseParagraphBreak

\secref{g2:sec:1501}先把流形放进高维空间，再寻找保留单射性与微分秩的投影方向。达到 \(2n+1\) 维时，维数比较足以避开所有割线方向；再降一维，一般位置会留下孤立双点。\secref{g2:sec:1502}专门处理这一层障碍。之后转向一个已经嵌入的子流形：\secref{g2:sec:1503}将其邻域与法丛联系起来，\secref{g2:sec:1504}给出边界的乘积邻域。这些邻域构造又使\secref{g2:sec:1505}能够延拓局部规定的运动。

\ProseParagraphBreak

最后考虑一个光滑参数族 \(p:E\to B\)。淹没条件说明每个点附近都有投影坐标，却未保证整条纤维能同时装进同一个乘积图。Ehresmann\footnote{Charles Ehresmann（1905—1979），法国数学家，研究微分拓扑、纤维丛与范畴论。}定理用固有性补上这一环：紧参数范围内的所有提升轨道留在一个紧集里，因而能把一条纤维完整地搬到相邻纤维。读到这里时，应特别留意“每一点的解短时间存在”与“全部初值都能走完指定时间”之间的差别。

\ProseParagraphBreak

本章默认流形有限维、Hausdorff、第二可数，未特别说明边界时均指无边界流形。前章的 Sard 定理、参数横截性与相对扰动用于选取一般位置；第一卷第11—13章的逆函数定理、单位分解、常微分方程光滑依赖与逃逸判据用于邻域和流。用到这些结论时会恢复具体条件。弱嵌入与流的参考见\cite[第6、9章]{Lee2012Smooth}；强 Whitney 界的双点构造参见\cite{Whitney1944SelfIntersections}。

% !TeX root = ../../Book2.tex
% 纲要标题：### 15.1 高维初始嵌入与泛型投影
% 纲要位置：几何学计划纲要.md，第 1576 行。
% 正文状态：初稿；尚未正式审核。

% 纲要内容开始
% 
% * 由局部坐标构造整体嵌入
% * 投影方向的坏集合
% * 弱 Whitney 维数界
% * 非紧源的穷竭及固有性
% 
% 纲要内容结束

\section{高维初始嵌入与泛型投影}
\label{g2:sec:1501}

一张坐标图把流形的一小块放进 Euclid 空间，但不同图中的坐标不能直接首尾相接。先让每张图负责几个独立的坐标分量，便能把局部信息全部保存下来；得到一个维数可能很高的嵌入后，再寻找不会压坏切向量、也不会合并不同点的投影方向。紧流形上的证明由这两步完成。非紧流形还须控制无穷远，否则单射浸入的像仍可能发生拓扑上的聚集。

\subsection{由局部坐标构造整体嵌入}
\label{g2:subsec:1501:01}

先回忆第一卷第11章第4节的区别：浸入要求微分单射；嵌入还要求映射是到子空间像的同胚。对紧源，连续单射到 Hausdorff 空间自动是到像的同胚，因此只需构造单射浸入。

\begin{proposition}\label{g2:prop:finite-coordinate-embedding}
	每个紧光滑 \(n\) 维流形 \(M\)，允许有边界，都能光滑嵌入某个有限维 Euclid 空间。
\end{proposition}
\begin{proof}
	由第一卷第12章第3节的相对紧坐标细化与截断构造，取有限张图 \((U_i,x_i)\)，以及较小开集 \(V_i\) 和光滑函数 \(\rho_i:M\to[0,1]\)，使 \(V_i\) 覆盖 \(M\)，\(\rho_i=1\) 于 \(\overline V_i\) 的邻域，\(\operatorname{supp}\rho_i\Subset U_i\)。把 \(\rho_i x_i\) 在 \(U_i\) 外补零，定义
	\[
		F(p)=\bigl(\rho_1(p),\rho_1(p)x_1(p),\ldots,
		\rho_s(p),\rho_s(p)x_s(p)\bigr)\in\mathbb R^{s(n+1)}.
	\]
	补零光滑，因为每个支集离相应图域的边缘留有邻域。

	若 \(F(p)=F(q)\)，选 \(i\) 使 \(p\in V_i\)。则 \(\rho_i(p)=\rho_i(q)=1\)，尤其 \(q\in U_i\)，而对应的向量块给出 \(x_i(p)=x_i(q)\)，故 \(p=q\)。在 \(p\in V_i\) 附近，第 \(i\) 个向量块就是 \(x_i\)，其微分单射，故 \(\mathrm{d}F_p\) 单射。最后，紧性保证连续逆，因而 \(F\) 为嵌入。边界图中的完整坐标微分同样可逆，所以证明也包括带边界情形。
\end{proof}

保留额外的标量 \(\rho_i\) 有实际作用：它告诉我们哪个坐标块正在工作，并允许从乘积 \(\rho_i x_i\) 恢复坐标。若只把加权坐标相加到同一个 \(\mathbb R^n\) 中，来自不同图的信息可能相互抵消，也可能把不同点混在一起。

\subsection{投影方向的坏集合}
\label{g2:subsec:1501:02}

设已有单射浸入 \(F:M^n\to\mathbb R^q\)。对单位向量 \(v\)，记 \(\pi_v:\mathbb R^q\to v^\perp\) 为正交投影。投影可能以两种方式失败：某条非零切向量被压成零，或者连接两个像点的割线被压成零。这两种情况都能把“坏方向”写成一个光滑映射的像。

\ProseParagraphBreak

在 \(M\) 上选一个辅助丛度量，记单位切丛为 \(U\mathrm{T}M\)。第一卷第12章第3节的丛度量存在定理保证这种选择；单位球在每个局部标架中是光滑超曲面，故 \(\dim U\mathrm{T}M=2n-1\)。考虑
\[
	\tau:U\mathrm{T}M\longrightarrow S^{q-1},\qquad
	\tau(x,u)=\frac{\mathrm{d}F_x(u)}{|\mathrm{d}F_x(u)|},
\]
以及去掉对角线的点对空间上的映射
\[
	\sigma:(M\times M)\setminus\Delta_M\longrightarrow S^{q-1},\qquad
	\sigma(x,y)=\frac{F(x)-F(y)}{|F(x)-F(y)|}.
\]
分母分别由浸入性和单射性保证非零。\(v\) 属于 \(\tau\) 的像，恰好意味着某个非零切向量落入 \(\ker\pi_v\)；\(v\) 属于 \(\sigma\) 的像，恰好意味着某对不同点被投影合并。由于 \(u\) 与 \(-u\)、\((x,y)\) 与 \((y,x)\) 都出现，这两个像已同时包含正、反方向。

\begin{proposition}\label{g2:prop:good-projection-directions}
	设 \(F:M^n\to\mathbb R^q\) 是光滑单射浸入，\(n\ge1\)，\(M\) 允许有边界。若 \(q>2n\)，则几乎每个单位方向 \(v\) 使 \(\pi_v\circ F\) 为浸入；若 \(q>2n+1\)，则几乎每个 \(v\) 同时使它单射。好方向在单位球中稠密。
\end{proposition}
\begin{proof}
	当 \(q>2n\) 时，\(\dim U\mathrm{T}M=2n-1<q-1\)，由\cref{g2:cor:sard-lower-dimension}，\(\tau\) 的像零测。当 \(q>2n+1\) 时，点对空间维数 \(2n<q-1\)，\(\sigma\) 的像也零测；两者的并仍零测。任一非空球面开集有正测度，所以补集稠密。

	若源有边界，将单位切丛分为底点在内部与边界的两部分，维数分别为 \(2n-1\)、\(2n-2\)。点对空间则分为内部乘内部、内部乘边界、边界乘内部、边界乘边界四个无边界流形，维数分别至多 \(2n,2n-1,2n-1,2n-2\)。在每一部分应用同一低维像结论，再取有限并即可。\(n=0\) 时没有切向条件；点对空间可数，割线方向的像也可数。
\end{proof}

\subsection{弱 Whitney 维数界}
\label{g2:subsec:1501:03}

\begin{theorem}[弱 Whitney 嵌入定理]\label{g2:thm:weak-whitney-compact}
	设 \(M\) 是紧光滑 \(n\) 维流形，允许有边界。则 \(M\) 可光滑嵌入 \(\mathbb R^{2n+1}\)。若 \(n\ge1\)，还存在光滑浸入 \(M\to\mathbb R^{2n}\)。
\end{theorem}
\begin{proof}
	若 \(n=0\)，紧零维流形是有限点集，直接选直线中的相同数目互异点即可。以下设 \(n\ge1\)。先由\cref{g2:prop:finite-coordinate-embedding}得到某个 \(F:M\to\mathbb R^q\)。若 \(q\le2n+1\)，补零坐标即可；否则用\cref{g2:prop:good-projection-directions}选好方向，将维数降一，所得仍是单射浸入。紧性使它仍为嵌入，所以可以重复，直到环境维数为 \(2n+1\)。从这里再投影一次，切向坏方向的维数仍小于目标方向球的维数，因此能保留浸入性，得到 \(\mathbb R^{2n}\) 中的浸入。
\end{proof}

最后一次投影有一个明确变化：割线映射的源与方向球同为 \(2n\) 维，零测论证已经不能排除其像覆盖一块开区域。于是微分仍可单射，不同点却可能投到一起。下一节的强界正要处理这些双点。

\begin{example}\label{g2:exm:circle-bad-projection}
	单位圆 \(S^1\subset\mathbb R^2\) 沿竖直方向投到横轴，成为函数 \((x,y)\mapsto x\)。圆的最高点和最低点有相同的像，而最左点和最右点的切线竖直，在那里投影微分为零。对圆周投到直线，转动方向只会转动这些坏点；维数条件不足时，没有理由期待避开全部坏方向。
\end{example}

\subsection{非紧源的穷竭及固有性}
\label{g2:subsec:1501:04}

现在令 \(M\) 为无边界光滑流形，仍假设有限维、Hausdorff 且第二可数。单射浸入未必是嵌入；第一卷第11章第4节的固有单射浸入命题给出可用的判据：若紧集的原像都紧，则到像的逆连续。下面会在构造过程中一直保留这个条件。

\begin{lemma}\label{g2:lem:smooth-proper-exhaustion}
	每个这样的 \(M\) 都有非负光滑函数 \(h:M\to[0,\infty)\)，使每个次水平集 \(h^{-1}([0,c])\) 紧。
\end{lemma}
\begin{proof}
	取第一卷第12章第3节构造中的紧穷竭 \(K_j\subset\operatorname{int}K_{j+1}\)，\(\bigcup_jK_j=M\)。以光滑截断选 \(\chi_j:M\to[0,1]\)，在 \(K_j\) 附近等于1，支集包含于 \(\operatorname{int}K_{j+1}\)。定义
	\[
		h=\sum_{j\ge1}(1-\chi_j).
	\]
	任一点附近属于某个 \(K_m\) 的内部，当 \(j\ge m\) 时对应项恒为零，所以该和局部有限且光滑。若 \(x\notin K_{m+1}\)，则前 \(m\) 项均为1，故 \(h(x)\ge m\)。对给定 \(c\)，选整数 \(m>c\)，次水平集便是 \(K_{m+1}\) 的闭子集，因此紧。
\end{proof}

为了得到有限维的初始嵌入，把 \(h\) 的相邻壳层分配到两个坐标块。由\cref{g2:thm:sard}，对每个整数 \(j\ge0\) 可选正则值
\[
	j<a_j<b_j<j+1.
\]
令 \(D_0=h^{-1}([0,1])\)、\(E_0=h^{-1}([0,a_1])\)，并对 \(j\ge1\) 令
\[
	D_j=h^{-1}([j,j+1]),\qquad E_j=h^{-1}([b_{j-1},a_{j+1}]).
\]
这些 \(E_j\) 是紧带边界流形，\(D_j\subset\operatorname{Int}E_j\)，且同为偶数或同为奇数的不同 \(E_j\) 两两不交。正则值条件在边界附近把 \(h\) 变成一个坐标，因此 \(E_j\) 的边界确为相应水平集，不会出现尖角。

\ProseParagraphBreak

由紧源情形，取嵌入 \(f_j:E_j\to\mathbb R^{2n+1}\)，以及 \(\beta_j\) 在 \(D_j\) 附近等于1、支集包含于 \(\operatorname{Int}E_j\) 的光滑截断。补零后定义
\[
	F(x)=\left(\sum_{j\text{ 偶}}\beta_j(x)f_j(x),\quad
	\sum_{j\text{ 奇}}\beta_j(x)f_j(x),\quad h(x)\right)
	\in\mathbb R^{4n+3}.
\]
每个坐标块在一点附近至多出现一项，且壳层族局部有限，所以 \(F\) 光滑。若 \(F(x)=F(y)\)，先由最后一坐标知 \(h(x)=h(y)\)；选 \(j\) 使 \(x\in D_j\)，便有 \(y\in D_j\)。相应奇偶块给出 \(f_j(x)=f_j(y)\)，从而 \(x=y\)。在 \(x\in D_j\) 附近，这一块就是 \(f_j\)，其微分单射，所以 \(F\) 为浸入。最后坐标 \(h\) 保证 \(F\) 固有：紧目标集的原像是某个紧次水平集的闭子集。因此 \(F\) 已经是有限维空间中的固有嵌入。

\ProseParagraphBreak

普通投影未必保留固有性。为控制它，先把前 \(4n+2\) 个坐标用微分同胚
\[
	z\longmapsto\frac{z}{\sqrt{1+|z|^2}}
\]
压进单位球，保留最后坐标 \(h\)。所得固有嵌入的像包含于某条直线 \(L\) 周围的有界半径管中。若再选投影方向 \(v\notin L\)，这一“管轴”仍留下非零方向，不能全部被压掉。

\begin{lemma}\label{g2:lem:proper-projection-tube}
	设 \(F:M\to\mathbb R^q\) 为固有映射，其像包含于直线 \(L\) 周围的半径 \(R\) 管中。若单位向量 \(v\notin L\)，则 \(\pi_v\circ F\) 固有，且其像仍包含于直线 \(\pi_v(L)\) 周围某个有界半径管中。
\end{lemma}
\begin{proof}
	写 \(z=s+b\)，其中 \(s\in L\)、\(|b|\le R\)。因为 \(v\notin L\)，\(\pi_v|_L\) 单射，故存在 \(c>0\)，使 \(|\pi_v(s)|\ge c|s|\)。若 \(|\pi_v(z)|\le B\)，则
	\[
		c|s|\le|\pi_v(z)|+|\pi_v(b)|\le B+R,
	\]
	从而 \(|z|\le(B+R)/c+R\)。因此紧目标集的原像闭且包含于 \(F\) 的一个紧集原像中，因而紧。另一方面，\(\pi_v(z)\) 到 \(\pi_v(L)\) 的距离不超过 \(|\pi_v(b)|\le R\)，得到管的控制。
\end{proof}

\begin{theorem}[弱 Whitney 嵌入定理的固有形式]\label{g2:thm:weak-whitney-proper}
	每个有限维、Hausdorff、第二可数的无边界光滑 \(n\) 维流形都能固有地光滑嵌入 \(\mathbb R^{2n+1}\)。若 \(n\ge1\)，还存在到 \(\mathbb R^{2n}\) 的固有光滑浸入。
\end{theorem}
\begin{proof}
	若 \(n=0\)，第二可数的零维流形是至多可数的离散集，将各点依次送到不同非负整数便是固有嵌入。以下设 \(n\ge1\)，使用上面构造的有界管内固有嵌入。当环境维数 \(q>2n+1\) 时，先避开\cref{g2:prop:good-projection-directions}的零测坏方向，再避开管轴的两个单位方向。这样得到的投影仍单射、浸入且固有，因而仍是嵌入；\cref{g2:lem:proper-projection-tube}还保留了下一步所需的管。迭代到 \(q=2n+1\) 即得第一项。若 \(n\ge1\)，再选一次避开切向坏方向和管轴的投影，保留浸入性与固有性，得到第二项。
\end{proof}

以上壳层与管轴构造也见\cite[第132--135页，引理6.13--6.14及定理6.15]{Lee2012Smooth}。紧性和固有性在证明中承担了不同工作：前者把有限张坐标图拼成初始嵌入，后者阻止非紧源在投影后从无穷远重新聚到有限位置。关于带边界非紧流形，可在\secref{g2:sec:1504}构造双化以后，将双化流形的固有嵌入限制到原来的闭半边。

% !TeX root = ../../Book2.tex
% 纲要标题：### 15.2 强 Whitney 界的附加论证
% 纲要位置：几何学计划纲要.md，第 1583 行。
% 正文状态：未撰写；本文件目前是框架。

% 纲要内容开始
% 
% * 一般位置的自交模型
% * 消去双点所需的局部构造
% * 高维和低维特例的分开处理
% * 不能由一次泛型投影得到最优界
% 
% 纲要内容结束

\section{强 Whitney 界的附加论证}
\label{g2:sec:1502}

待撰写

\subsection{一般位置的自交模型}
\label{g2:subsec:1502:01}
待撰写

\subsection{消去双点所需的局部构造}
\label{g2:subsec:1502:02}
待撰写

\subsection{高维和低维特例的分开处理}
\label{g2:subsec:1502:03}
待撰写

\subsection{不能由一次泛型投影得到最优界}
\label{g2:subsec:1502:04}
待撰写
% 纲要条目排为编号小节标题，并进入目录与书签；说明段仍保留于纲要注释中。

% !TeX root = ../../Book2.tex
% 纲要标题：### 15.3 法丛与管状邻域
% 纲要位置：几何学计划纲要.md，第 1590 行。
% 正文状态：初稿；尚未正式审核。

% 纲要内容开始
% 
% * 法向指数映射
% * 零截面附近的逆函数论证
% * 非紧子流形的变半径控制
% * 管状邻域及其选择无关性
% 
% 纲要内容结束

\section{法丛与管状邻域}
\label{g2:sec:1503}

在子流形的一点附近，适配坐标把它写成 \(\mathbb R^k\times\{0\}\)。沿整个子流形移动时，这些横向坐标可能发生扭转，所以邻域通常不能写成一个固定乘积。法丛恰好记录横向方向怎样拼接。管状邻域定理说明：只要把每条纤维取得足够短，这个线性模型就能实际描述子流形的一个开邻域。

\subsection{法向指数映射}
\label{g2:subsec:1503:01}

设 \(S\subset N\) 是无边界流形 \(N\) 的无边界嵌入子流形。\secref{g2:sec:1403}已定义内在的法丛
\[
	\nu_N S=\mathrm{T}N|_S/\mathrm{T}S.
\]
一个法向量首先是模去切向量后的等价类。由第一卷第12章第3节的丛度量存在及正交补命题，在 \(N\) 上选光滑 Riemann 度量 \(g\)，便能把每个等价类 \(v\in\nu_NS\) 唯一地表示为正交向量 \(\iota_g(v)\in(\mathrm{T}S)^\perp\)。正交投影的局部矩阵由光滑正定矩阵的逆给出，所以 \(\iota_g\) 随基点光滑。

\ProseParagraphBreak

还要把这个向量变成附近的一个实际点。为此，在 \(\mathrm{T}N\) 上选择一个光滑联络 \(\nabla\)；第一卷第15章第1节用局部平凡联络的单位分解加权证明了其存在。这里只需要联络给出的短时曲线，度量负责选出法向代表。若坐标中的联络系数为 \(\Gamma^a_{bc}\)，初值为 \((x,v)\) 的曲线满足
\[
	\ddot z^a+\Gamma^a_{bc}(z)\dot z^b\dot z^c=0,
	\qquad z(0)=x,\quad\dot z(0)=v.
\]
把 \((z,\dot z)\) 当作一阶方程的未知量，第一卷第13章第1节的存在唯一性与初值光滑依赖定理适用。\(v=0\) 时解是常曲线，因此对 \((x,0)\) 的充分小邻域，解都能存在到时间1。令
\[
	\exp_x^\nabla(v)=z_{x,v}(1).
\]
它在 \(\mathrm{T}N\) 的零截面某邻域上光滑，称为该联络的局部指数映射。方程由 \(\nabla_{\dot z}\dot z=0\) 给出，与坐标选择相容。

\ProseParagraphBreak

在常解 \((x,0)\) 处，对初速度作一阶变化，二次速度项的一阶变化为零，故变分方程为 \(J''=0\)、\(J(0)=0\)、\(J'(0)=w\)。于是 \(J(1)=w\)，即
\[
	\mathrm{d}_v\exp_x^\nabla\big|_{v=0}=\operatorname{id}_{\mathrm{T}_xN}.
\]
同时 \(\exp_x^\nabla(0)=x\)，所以在基点方向上的微分也是恒等。定义法向映射
\[
	\Phi(v)=\exp_x^\nabla(\iota_gv),\qquad v\in\nu_xS.
\]
选择任意联络已足够；此处没有要求测地线在任意长时间存在，也没有使用最短路径性质。

\subsection{零截面附近的逆函数论证}
\label{g2:subsec:1503:02}

零截面处的切空间有规范分解
\[
	\mathrm{T}_{0_x}(\nu_NS)\cong\mathrm{T}_xS\oplus\nu_xS.
\]
第一项来自零截面的微分，第二项是纤维中的速度。在局部丛坐标中，这就是基底坐标与纤维坐标的分解；零纤维坐标使过渡函数中对基点求导的附加项消失，因此分解与平凡化无关。

\begin{proposition}\label{g2:prop:normal-exponential-local}
	对上面构造的 \(\Phi\)，有 \(\Phi(0_x)=x\)，且
	\[
		\mathrm{d}\Phi_{0_x}(u,v)=u+\iota_g(v).
	\]
	故 \(\Phi\) 在每个 \(0_x\) 的某邻域内为局部微分同胚。
\end{proposition}
\begin{proof}
	基点方向和纤维方向的微分已分别算出。由于 \(\mathrm{T}_xN=\mathrm{T}_xS\oplus(\mathrm{T}_xS)^\perp\)，右边是线性同构。第一卷第11章第3节的逆函数定理遂给出局部逆。
\end{proof}

这个结论还没有把整个零截面附近的不同纤维分开。即使每一点附近都有逆，来自相距很远的两个基点的短向量仍可能被送到同一点。紧子流形可用紧性统一选择半径；非紧情形需要让半径随位置缩小。

\begin{example}\label{g2:exm:shrinking-circle-tubes}
	在平面中取圆周的不交并
	\[
		S=\coprod_{j\ge1}\{(3j,0)+j^{-1}(\cos\theta,\sin\theta):\theta\in\mathbb R\}.
	\]
	这是闭的非紧嵌入一维子流形。采用 Euclid 度量和直线指数时，若统一法向半径为 \(\varepsilon>0\)，选 \(j\) 使 \(j^{-1}<\varepsilon\)，则第 \(j\) 个圆的所有向心法向量都可到达圆心，法向映射不再单射。各圆上的半径改为小于 \(1/(2j)\)，则这些问题消失。局部逆的存在并不保证统一的正半径。
\end{example}

\subsection{非紧子流形的变半径控制}
\label{g2:subsec:1503:03}

以下缩小论证既用于法向指数，也会用于比较不同的管状映射。保留紧参数族的版本，可以同时处理一段同位，而不用对无穷多个参数分别选邻域。

\begin{lemma}\label{g2:lem:zero-section-shrinking}
	设 \(S\subset N\) 是闭的无边界嵌入子流形，\(E\to S\) 为有丛度量的光滑向量丛，\(\dim E=\dim N\)。设 \(\Phi_t\) 是依赖于 \(t\in[0,1]\) 的联合光滑映射，在所有零向量附近有定义，并满足
	\[
		\Phi_t(0_x)=x,\qquad \mathrm{d}(\Phi_t)_{0_x}\text{ 为同构}.
	\]
	则存在连续函数 \(\varepsilon:S\to(0,\infty)\)，使每个 \(\Phi_t\) 在
	\[
		E_\varepsilon=\{v\in E:|v|<\varepsilon(\pi(v))\}
	\]
	上都是到 \(N\) 中开集的微分同胚；各局部逆对 \(t\) 也光滑。只有一个映射时，视为常参数族即可。
\end{lemma}
\begin{proof}
	由\cref{g2:lem:smooth-proper-exhaustion}，选非负光滑固有函数 \(h:N\to[0,\infty)\)，记 \(K_j=S\cap\{h\le j\}\)。闭性保证 \(K_j\) 紧。对每个 \(j\)，存在 \(r_j>0\)，使所有 \(t\) 下、基点在 \(K_j\) 上且长度小于 \(r_j\) 的向量都在定义域内，微分可逆，而且
	\[
		|h(\Phi_t(v))-h(\pi(v))|<1/3.
	\]
	这是对紧集 \([0,1]\times\{0_x:x\in K_j\}\) 使用有限覆盖和连续性。还可要求同一个 \(t\) 下，这些向量之间没有重合的像。若做不到，便有 \(t_i\)、不同的 \(v_i,w_i\)，基点都在 \(K_j\)，长度都趋零，且 \(\Phi_{t_i}(v_i)=\Phi_{t_i}(w_i)\)。取子列使 \(t_i\to t\)、\(\pi(v_i)\to x\)、\(\pi(w_i)\to y\)。极限等式给出 \(x=y\)。但映射 \((t,v)\mapsto(t,\Phi_t(v))\) 在 \((t,0_x)\) 的微分可逆，故其局部单射性排除了这样的序列。缩小后可令 \(r_j\) 单调不增。

	选连续正函数 \(\varepsilon(x)\)，使当 \(j-1\le h(x)\le j\) 时有 \(\varepsilon(x)<r_{j+3}\)。例如先在整数点选单调递减且足够小的正数，作分段线性插值，再与 \(h|_S\) 复合即可。必要时用\cref{g2:thm:relative-smooth-approximation}将这个正函数换成更小的光滑正函数。

	设 \(v,w\in E_\varepsilon\)，且 \(\Phi_t(v)=\Phi_t(w)\)。上述高度估计给出 \(|h(\pi(v))-h(\pi(w))|<2/3\)。若 \(j-1\le h(\pi(v))\le j\)，则两个基点都属于 \(K_{j+1}\)，而两个向量的长度都小于 \(r_{j+1}\)：第二个基点至多比第一个提前或推后一层，半径选择中预留的三层已涵盖这一差别。\(K_{j+1}\) 上的单射性于是给出 \(v=w\)。故每个 \(\Phi_t\) 在整个 \(E_\varepsilon\) 上单射。

	它的微分处处可逆，因此是开映射；单射的局部微分同胚便是到其开像的微分同胚。参数逆的光滑性同样来自 \((t,v)\mapsto(t,\Phi_t(v))\) 的局部逆，各局部逆由唯一性相容。
\end{proof}

上面的紧集论证还允许法丛随参数变化。具体地，设 \(K\) 紧且无边界，\(i_t:K\to N\) 为光滑嵌入族，\(E\to[0,1]\times K\) 是光滑向量丛。若一个零截面附近的光滑映射
\[
	\mathcal F(t,x,v)=(t,F(t,x,v))
\]
满足 \(F(t,x,0)=i_t(x)\)，且 \(\mathrm{d}\mathcal F\) 在零截面处可逆，则它在某个共同的小圆盘丛上仍是到相对开像的微分同胚。事实上，带时间的零截面像由紧空间 \([0,1]\times K\) 嵌入 \([0,1]\times N\)；若任意小半径仍有两个向量像重合，取趋零向量及基点的收敛子列，极限像相等便迫使时间和基点都相等，再由该点的局部逆排除重合。零截面紧同时给出定义域与可逆微分的统一半径。此论证只用丛的局部平凡化，不要求把所有 \(E_{t,x}\) 预先识别成一个固定丛。在参数端点，先局部延拓系数；由于第一坐标始终是 \(t\)，逆函数定理所得的局部逆也保持时间，因而可限制回半空间。这一版本将在同位延拓中保持每个时刻的子流形。

若 \(S\) 原来没有在 \(N\) 中闭，先取所有适配坐标域的并 \(N_0\)。它是包含 \(S\) 的开集，且 \(S\) 在 \(N_0\) 中闭：任何 \(z\in N_0\) 都位于某张适配图内，图中的坐标切片是相对闭集，因而一个不在切片上的点不可能是 \(S\) 的聚点。在 \(N_0\) 中应用引理即可。由此，管状邻域存在性不要求原嵌入固有，但所需半径可以趋于零。

\subsection{管状邻域及其选择无关性}
\label{g2:subsec:1503:04}

\begin{definition}\label{g2:def:tubular-neighborhood}
	设 \(S\subset N\) 为无边界嵌入子流形。一张\term[guanzhuanglinyu]{管状邻域映射}[tubular neighborhood map]是从 \(\nu_NS\) 的零截面某邻域 \(V\) 到 \(N\) 的开集 \(U\supset S\) 的微分同胚 \(\tau:V\to U\)，满足 \(\tau(0_x)=x\)，并且纤维方向的微分在取法商后为恒等：
	\[
		\nu_xS\longrightarrow\mathrm{T}_xN\longrightarrow\nu_xS,
		\qquad v\longmapsto\mathrm{d}\tau_{0_x}(0,v)\longmapsto v.
	\]
	称 \(U\) 为 \(S\) 的\term[guanzhuanglinyu]{管状邻域}[tubular neighborhood]。可缩小 \(V\)，使它为上一引理中的变半径圆盘丛。
\end{definition}

\begin{theorem}[管状邻域定理]\label{g2:thm:tubular-neighborhood}
	每个无边界光滑流形中的无边界嵌入子流形都有管状邻域。子流形紧时可取常数法向半径，一般情形可取正的变半径。
\end{theorem}
\begin{proof}
	度量与联络给出的法向指数 \(\Phi\) 满足\cref{g2:prop:normal-exponential-local}，故\cref{g2:lem:zero-section-shrinking}在适当环境开集中给出 \(\tau=\Phi|_{E_\varepsilon}\)。其零截面值正确，纤维微分是 \(\iota_g\)，取商后为恒等，因此满足定义。若 \(S\) 紧，连续正函数 \(\varepsilon\) 有正的下界，进一步取更小常数即可。
\end{proof}

例如 \(S^k\subset\mathbb R^{k+1}\) 的外法向线丛平凡，\(\tau(x,t)=(1+t)x\)、\(|t|<1/2\)，给出一条环形邻域。在 Möbius 带总空间中，零截面的法线丛仍是 Möbius 线丛；其邻域是扭转的带，而非圆周与区间的乘积。管状模型保留的是整个法丛的拼接。

\ProseParagraphBreak

在 \(U=\tau(E_\varepsilon)\) 上，令 \(r=\pi\circ\tau^{-1}\)，得到光滑淹没 \(r:U\to S\)，且 \(r|_S=\operatorname{id}\)。纤维缩放给出
\[
	R_s(\tau(v))=\tau((1-s)v),\qquad0\le s\le1,
\]
因此 \(U\) 光滑形变收缩到 \(S\)。变半径条件只依赖基点，故整条线段 \((1-s)v\) 都留在定义域内。这些结论常用来把邻域中的问题送回子流形。\cref{g2:fig:tubular-retraction}画出圆周的情形。

\begin{figure}[htbp]
	\centering
	% !TeX root = ../Book2.tex
\begin{tikzpicture}[scale=1.25,>=stealth]
	\fill[blue!8,even odd rule] (0,0) circle (1.35) (0,0) circle (0.65);
	\draw[dashed,blue!35] (0,0) circle (1.35);
	\draw[dashed,blue!35] (0,0) circle (0.65);
	\draw[thick,blue!65!black] (0,0) circle (1);
	\foreach \a in {0,45,90,135,180,225,270,315} {
		\draw[gray] ({0.68*cos(\a)},{0.68*sin(\a)})--({1.32*cos(\a)},{1.32*sin(\a)});
	}
	\draw[->,thick] (1,0)--(1.3,0) node[right] {\(v\)};
	\fill (1,0) circle (1pt) node[below left] {\(x\)};
	\node at (-0.65,1.35) {\(S\)};
	\draw[->] (1.65,0.55)--(3.15,0.55) node[midway,above] {\(r\)};
	\node[align=center] at (2.4,-0.2) {沿法向纤维\\回到零截面};
	\draw[thick,blue!65!black] (4.5,0) circle (1);
	\fill (5.5,0) circle (1pt) node[right] {\(x\)};
\end{tikzpicture}

	\caption[法向纤维与管状邻域收缩]{圆周附近的短法向纤维填成环形邻域。投影 \(r\) 将同一纤维上的点送回基点 \(x\)；纤维缩放则给出从恒等映射到 \(r\) 的形变。}
	\label{g2:fig:tubular-retraction}
\end{figure}


\begin{proposition}\label{g2:prop:tubular-uniqueness}
	设 \(\tau_0,\tau_1\) 是同一个无边界嵌入子流形 \(S\subset N\) 的两张管状邻域映射，均按上述定义正规化法方向微分。存在法丛零截面的一个共同邻域及联合光滑族 \(\tau_t\)、\(0\le t\le1\)，使每张 \(\tau_t\) 都是管状邻域映射，端点为原映射的限制，并且整族在零截面上不变。
\end{proposition}
\begin{proof}
	先缩小到 \(\tau_1(v)\) 落在 \(\tau_0\) 的像内，令 \(F=\tau_0^{-1}\circ\tau_1\)。它是法丛零截面附近的微分同胚，固定零截面，且诱导的法方向微分为恒等。记纤维缩放为 \(\delta_t(v)=tv\)。对 \(t>0\)，定义
	\[
		F_t=\delta_{1/t}\circ F\circ\delta_t.
	\]
	在局部丛坐标中写 \(F(x,v)=(a(x,v),b(x,v))\)。零截面与正规化条件给出
	\[
		a(x,0)=x,\qquad b(x,0)=0,\qquad \partial_vb(x,0)=I.
	\]
	于是
	\[
		F_t(x,v)=\left(a(x,tv),\frac{b(x,tv)}t\right),\qquad
		\frac{b(x,tv)}t=\int_0^1\partial_vb(x,stv)v\,\mathrm{d}s.
	\]
	积分表达说明这族光滑延到 \(t=0\)，且 \(F_0(x,v)=(x,v)\)。这些表达来自整体纤维缩放，故不同平凡化中的延拓相容。零截面处的微分形如
	\[
		\mathrm{d}(F_t)_{0_x}(u,v)=(u+tB_xv,v),
	\]
	始终可逆。

	令 \(\tau_t=\tau_0\circ F_t\)。紧参数区间保证它们在零截面的某共同邻域都有定义；对这一族使用\cref{g2:lem:zero-section-shrinking}，取得共同的变半径邻域，使各张都为到开像的微分同胚。它们固定零截面，法方向微分都是恒等，且 \(t=0,1\) 时分别为 \(\tau_0,\tau_1\)。
\end{proof}

这就是这里的选择无关性：缩小邻域后，度量、联络和半径所给出的不同管状坐标可以通过固定子流形的光滑同位相连。它并未指定唯一的一张坐标图。若一个邻域参数化在法商上诱导的是丛自同构 \(A\)，先把纤维变量换成 \(A^{-1}v\)，便得到定义中的正规化；把这项识别写清楚，才能比较相同法丛上的两张映射。Euclid 空间中的法向加法与变半径构造还可参看\cite[第138--141页，定理6.24及命题6.25]{Lee2012Smooth}。

% !TeX root = ../../Book2.tex
% 纲要标题：### 15.4 边界的领状邻域
% 纲要位置：几何学计划纲要.md，第 1597 行。
% 正文状态：初稿；尚未正式审核。

% 纲要内容开始
% 
% * 向内横截向量场与局部流
% * 边界邻域的整体拼接
% * 双化和相对粘合中的应用
% 
% 纲要内容结束

\section{边界的领状邻域}
\label{g2:sec:1504}

在半空间图中，边界附近看起来像“边界坐标加一个非负坐标”。各张图的非负坐标并不相同，因而还不能把所有边界点统一向内推一段距离。本节用一个向内向量场把这些局部图连接起来：每条轨道从边界出发，时间成为共同的横向坐标。由此得到的领状邻域可用于双化、粘合和保持边界数据的构造。

\subsection{向内横截向量场与局部流}
\label{g2:subsec:1504:01}

设 \(M\) 为有限维、Hausdorff、第二可数的光滑带边界流形。在边界坐标 \((y,r)\)、\(r\ge0\) 中，向量 \(v\in\mathrm{T}_xM\) 严格指向内部的条件是 \(\mathrm{d}r_x(v)>0\)。第一卷第11章的边界坐标变换说明，另一边界定义坐标沿边界的微分是 \(\mathrm{d}r\) 的正倍数，所以这一条件不依赖所选图。它也说明同一点的向内向量组成凸集，正的加权平均仍向内。

\begin{lemma}\label{g2:lem:boundary-defining-function}
	每个光滑带边界流形 \(M\) 都有光滑函数 \(r:M\to[0,\infty)\)，满足 \(r^{-1}(0)=\partial M\) 且 \(\mathrm{d}r\ne0\) 于 \(\partial M\)。在边界的某个开邻域上，可选择光滑向量场 \(X\)，使 \(Xr=1\)。
\end{lemma}
\begin{proof}
	取由边界坐标图和内点坐标图组成的局部有限覆盖，并取从属光滑单位分解 \(\chi_i\)。边界图上记最后一个坐标为 \(r_i\)，内点图上令 \(r_i=1\)，定义 \(r=\sum_i\chi_i r_i\)，各项在支集外补零。每个内点的所有非零项都为正，故 \(r>0\)；在边界点，来自内点图的项在该点附近消失，其余 \(r_i\) 为零，所以 \(r=0\)。边界上
	\[
		\mathrm{d}r=\sum_i\chi_i\,\mathrm{d}r_i.
	\]
	各个 \(\mathrm{d}r_i\) 是同一个向内法向余向量的正倍数，且权重之和为1，故和非零。

	在各边界图中取 \(\partial_{r_i}\)，用单位分解加权拼成边界附近的光滑向量场 \(Y\)。其边界值严格向内，因此 \(Yr>0\) 在边界的某开邻域 \(U\) 上成立。令 \(X=Y/(Yr)\) 于 \(U\)，便有 \(Xr=1\)。这里使用的是第一卷第12章第3节已经证明的局部有限单位分解；局部有限性保证求和及逐项求导有意义。
\end{proof}

\(r\) 称为一个\term[bianjiedingyihanshu]{边界定义函数}[boundary defining function]。条件 \(\mathrm{d}r\ne0\) 保证它真正测量横向变化；仅有同一零点集的函数 \(r^2\) 不具有这个性质。归一化 \(Xr=1\) 的作用也很具体：沿从边界出发的积分曲线，\(r\) 的值恰好就是已经经过的时间。

\ProseParagraphBreak

在边界处不能直接要求 \(X\) 生成保持 \(M\) 的双侧流。向内的速度在负时间会走到半空间外。我们在边界坐标中把光滑分量局部延拓到 Euclid 开集，应用第一卷第13章第1节的初值存在、唯一和光滑依赖定理，再保留非负时间。因 \(r(0)=0\)、\(r'=1\)，所得小段正时间曲线位于内部。不同局部延拓在半空间内给出同一个方程，解的唯一性使这些正时间解一致。记其为
\[
	F(x,t)=\Phi_t(x),\qquad x\in\partial M,\quad 0\le t<\varepsilon_x.
\]
在 \(t=0\) 处，\(\mathrm{d}F\) 把边界切向量保持为原向量，并把 \(\partial_t\) 送到 \(X_x\)。这个向量不切于边界，因此上述导数为同构，局部逆函数论证给出边界附近的半空间坐标。

\subsection{边界邻域的整体拼接}
\label{g2:subsec:1504:02}

\begin{definition}\label{g2:def:collar}
	设 \(M\) 为光滑带边界流形。若光滑嵌入 \(c:\partial M\times[0,1)\to M\) 满足 \(c(x,0)=x\)，且像是 \(\partial M\) 在 \(M\) 中的开邻域，就称它为一个\term[lingzhuanglinyu]{领状邻域}[collar neighborhood]的参数化，简称领圈。
\end{definition}

\begin{theorem}[领状邻域定理]\label{g2:thm:collar-neighborhood}
	每个有限维、Hausdorff、第二可数的光滑带边界流形都存在领圈。边界可以非紧；领圈的像可取在任何预先给定的边界开邻域内。
\end{theorem}
\begin{proof}
	在前面构造的开集 \(U\) 上使用 \(Xr=1\) 的向量场，再把 \(U\) 与指定的边界邻域相交。由局部存在与光滑依赖，每个边界点有一个边界邻域 \(V_i\) 及正数 \(a_i\)，使从 \(V_i\) 出发的解在 \(0\le t<a_i\) 上存在并留在 \(U\)。可取 \(\{V_i\}\) 局部有限。设 \(\{\psi_i\}\) 为从属的单位分解，令
	\[
		\varepsilon(x)=\frac12\sum_i\psi_i(x)a_i>0.
	\]
	对每个 \(x\)，在 \(\psi_i(x)>0\) 的有限多个指标中选 \(a_i\) 最大者，则 \(\varepsilon(x)<a_i\) 且 \(x\in V_i\)。故 \(F\) 在
	\[
		D=\{(x,t):x\in\partial M,\ 0\le t<\varepsilon(x)\}
	\]
	上有定义，光滑且满足 \(r(F(x,t))=t\)。

	若 \(F(x,t)=F(y,s)\)，对两边取 \(r\) 得 \(t=s\)。两条轨道从相同时间、相同位置反向求解，在返回边界前都位于 \(M\)，唯一性给出 \(x=y\)。所以 \(F\) 在整个 \(D\) 上单射，这一步同时排除了来自遥远边界点的相撞。

	在正时间，流的微分把 \(\mathrm{T}_x\partial M\) 与 \(X_x\) 组成的基送到终点，并保持生成方向 \(X\)，故 \(\mathrm{d}F\) 仍可逆。也可注意 \(r\circ F=t\)，再在每张 \(t\) 切片上用流的局部逆。零时间处已由半空间局部逆得到同样结论。于是 \(F\) 是单射局部微分同胚，其局部逆在交叠处一致，拼成到开像的光滑整体逆。

	最后令 \(c(x,u)=F(x,u\varepsilon(x))\)，\(0\le u<1\)。\((x,u)\mapsto(x,u\varepsilon(x))\) 是到 \(D\) 的微分同胚，因而 \(c\) 是所需领圈。
\end{proof}

非紧边界上，轨道允许走的时间随点变化；最后的重参数化把这个变化吸收到 \(\varepsilon(x)\) 中。统一参数区间 \([0,1)\) 因而不表示每个点都能沿最初向量场运动同样长的时间。相关流与领圈构造可参看\cite[第217--223页，定理9.20、9.24、9.25]{Lee2012Smooth}。

\begin{example}\label{g2:exm:disk-collar}
	闭单位球 \(D^n\) 的一个领圈为 \(c(x,t)=(1-t/2)x\)，其中 \(x\in S^{n-1}\)、\(0\le t<1\)。其像是 \(\{z:1/2<|z|\le1\}\)，在 \(D^n\) 中开。边界点被逐渐推入内部，角方向保持不变；半径 \(1/2\) 没有包含在像中，对应参数端点 \(t=1\) 未取到。
\end{example}

\subsection{双化和相对粘合中的应用}
\label{g2:subsec:1504:03}

取 \(M\) 的两份副本 \(M_+,M_-\)，将对应的边界点识别，得到\term[shuanghua]{双化}[double]
\[
	\operatorname{Dbl}(M)=(M_+\amalg M_-)/\{x_+\sim x_-:x\in\partial M\}.
\]
给定领圈 \(c\) 后，用一个有符号的参数把接缝附近写成
\[
	\Theta:\partial M\times(-1,1)\longrightarrow\operatorname{Dbl}(M),\qquad
	\Theta(x,s)=\begin{cases}
		[c(x,s)]_+,&s\ge0,\\
		[c(x,-s)]_-,&s\le0.
	\end{cases}
\]
两条表达在零处一致。它是到接缝开邻域的同胚：在两半上分别有领圈的逆，它们的第一分量在边界一致，第二分量分别取正号、负号，因而由粘合定义得到连续逆。

\begin{proposition}\label{g2:prop:smooth-boundary-gluing}
	设 \(M_1,M_2\) 为同维光滑带边界流形，\(\varphi:\partial M_1\to\partial M_2\) 为微分同胚。选定两侧领圈后，商空间 \(M_1\cup_\varphi M_2\) 具有光滑无边界流形结构，使两份 \(M_i\) 都以闭的带边界子流形嵌入其中。特别地，\(\operatorname{Dbl}(M)\) 是无边界光滑流形。只沿边界分支的并作同样粘合时，其余边界保留。
\end{proposition}
\begin{proof}
	以 \(\partial M_1\) 为共同边界，在负半边用 \(c_2(\varphi(x),-s)\) 代替上面的表达。两侧内点保留原图；接缝附近取 \(\partial M_1\) 的坐标与 \(s\) 构成乘积图。两张接缝图的变换为 \((y,s)\mapsto(h(y),s)\)，所以光滑。接缝图与内点图的交落在 \(s>0\) 或 \(s<0\)，过渡由领圈、原流形坐标及 \(\varphi\) 的光滑性保证。

	商空间还须满足分离性与可数性。每个接缝点有上述乘积邻域，两侧内点也有原来的开邻域。两个不同接缝点用边界上的不交开集及较小领圈分离；内点与接缝点用避开该内点的小领圈分离；内点之间使用原流形的 Hausdorff 性或两侧互不相交的内部。这证明 Hausdorff 性。内部图及接缝的乘积图均可取可数族，故整体第二可数。各侧在这些坐标中为闭半空间，其补集是另一侧内部，所以是所述闭嵌入。
\end{proof}

粘合光滑对象时，领圈还告诉我们应匹配哪些数据。若两边的函数在领圈内都等于同一个边界函数 \(a(x)\)，它们就粘成光滑函数，因为接缝表达为与 \(s\) 无关的 \(a(x)\)。只要求边界值相等则不够：两侧都取领圈坐标 \(u\)，接起来变成 \(|s|\)，在零处不可微。对一般函数，需要各阶法向导数按 \(u=-s\) 的符号相容；先在领圈内调整为乘积形式，是一种直接保证相容的做法。\cite[第224--225页，定理9.29]{Lee2012Smooth}给出了相同的流形粘合构造。

% !TeX root = ../../Book2.tex
% 纲要标题：### 15.5 同位延拓
% 纲要位置：几何学计划纲要.md，第 1603 行。
% 正文状态：初稿；尚未正式审核。

% 纲要内容开始
% 
% * 同位的速度场
% * 管状邻域延拓与截断
% * 环境流及支集控制
% * 相对同位和非紧限制
% 
% 纲要内容结束

\section{同位延拓}
\label{g2:sec:1505}

把一个圆在平面中平移，可以同时平移整个平面；若只想移动圆附近的区域，还可以把平移速度截断。一般的嵌入族也提出同样的问题：已经规定子流形上的每个点怎样运动，能否让整个环境流形一起运动，并使远处保持原样？困难在于不同时间的速度给在不同位置上。先把这些速度写成随时间变化的向量场，再延拓到环境中，流的唯一性就能保证它准确实现原来的运动。

\subsection{同位的速度场}
\label{g2:subsec:1505:01}

\begin{definition}\label{g2:def:smooth-isotopy}
	设 \(K,N\) 为光滑流形。光滑映射 \(i:[0,1]\times K\to N\) 若对每个 \(t\) 都使 \(i_t=i(t,-)\) 成为嵌入，就称为一个光滑\term[tongwei]{同位}[isotopy]。若光滑映射 \(\Phi:[0,1]\times N\to N\) 的每个 \(\Phi_t\) 都是微分同胚，且 \(\Phi_0=\operatorname{id}_N\)，则称其为\term[huanjingtongwei]{环境同位}[ambient isotopy]。
\end{definition}

同位指定的是一族带参数的嵌入，而不只是一族集合。即使 \(i_t(K)\) 始终是同一个子集，其上的点也可能沿子流形移动。例如 \(i_t(e^{\mathrm i\theta})=e^{\mathrm i(\theta+t)}\) 保持圆周的像不变，却让每个标记点转动。环境同位若延拓这个族，应满足 \(\Phi_t\circ i_0=i_t\)，而非仅把两个像集合相互送到。

\ProseParagraphBreak

对光滑同位定义
\[
	v(t,i_t(x))=\frac{\partial}{\partial t}i(t,x)\in\mathrm{T}_{i_t(x)}N.
\]
因为 \(i_t\) 单射，右端唯一确定左端；这个向量属于环境切空间，未必切于 \(i_t(K)\)。为理解其光滑性，考虑带时间标记的映射
\[
	\widehat i:[0,1]\times K\longrightarrow[0,1]\times N,
	\qquad(t,x)\longmapsto(t,i_t(x)).
\]
其微分为 \((a,u)\mapsto(a,a\partial_ti+\mathrm{d}i_t(u))\)，所以单射。若 \(K\) 紧，\(\widehat i\) 又是从紧空间到 Hausdorff 空间的连续单射，因而是嵌入。于是在其像上，\(v\) 是向量丛 \(\operatorname{pr}_N^*\mathrm{T}N\) 的光滑截面。这正是需要延拓的数据：时间分量保持为零，空间分量等于指定速度。

\subsection{管状邻域延拓与截断}
\label{g2:subsec:1505:02}

先取紧无边界的 \(K\)。在 \(N\) 上选度量，令 \(E_{t,x}\) 为 \(\mathrm{d}i_t(\mathrm{T}_xK)\) 的正交补。这些空间在 \([0,1]\times K\) 上组成光滑向量丛：在局部标架中，正交投影由光滑的 Gram 矩阵及其逆给出。用\secref{g2:sec:1503}的法向指数构造
\[
	\mathcal E(t,x,u)=\bigl(t,\exp_{i_t(x)}u\bigr),\qquad u\in E_{t,x}.
\]
在零截面上它的微分可逆，且第一分量仍为 \(t\)。\cref{g2:lem:zero-section-shrinking}的紧参数版本给出零截面邻域，使 \(\mathcal E\) 成为到 \([0,1]\times N\) 中轨迹邻域的微分同胚。其逆的基点分量定义保持时间的光滑收缩
\[
	R(t,y)=(t,i_t(x)).
\]
这里的法向坐标只用于把空间点拉回同一时刻的子流形。

\ProseParagraphBreak

有了 \(R\)，向量场仍不能直接写成 \(v\circ R\)：这个向量位于 \(\mathrm{T}_{i_t(x)}N\)，而需要的值应位于 \(\mathrm{T}_yN\)。可以用局部标架明确解决这个问题。取覆盖轨迹的有限个向量丛平凡化，记环境切丛的局部标架为 \(e_{\alpha j}\)，并在轨迹上取从属单位分解 \(\lambda_\alpha\)。将
\[
	v=\sum_j v_\alpha^j e_{\alpha j}
\]
中的系数 \(v_\alpha^j\) 及权重 \(\lambda_\alpha\) 与 \(R\) 复合，再在 \((t,y)\) 处使用标架 \(e_{\alpha j}(y)\)。缩小邻域，使涉及的基点和 \(y\) 都落在该平凡化内，就得到
\[
	\widetilde X(t,y)=\sum_{\alpha,j}\lambda_\alpha(R(t,y))\,
	v_\alpha^j(R(t,y))\,e_{\alpha j}(y).
\]
在原轨迹上 \(R\) 是恒等，权重之和为1，故 \(\widetilde X=v\)。各项按其支集补零，得到光滑延拓。

\ProseParagraphBreak

设 \(\Omega\subset N\) 是包含整个运动轨迹 \(i([0,1]\times K)\) 的开集。轨迹紧，因此可在上述时空邻域中选一个光滑截断 \(\chi\)，于轨迹附近等于1，支集含于 \([0,1]\times\Omega\)，且支集的空间投影是紧集。第一卷第12章的光滑截断定理适用于这个带边界的时空区域。令 \(X_t(y)=\chi(t,y)\widetilde X(t,y)\)，在其余地方补零，则
\[
	X_t(i_t(x))=\partial_ti(t,x),\qquad
	\bigcup_{0\le t\le1}\operatorname{supp}X_t\subset L\Subset\Omega.
\]
一个共同的紧集 \(L\) 控制所有时刻的支集，是下一步得到整个时间区间上环境流的关键。

\subsection{环境流及支集控制}
\label{g2:subsec:1505:03}

\begin{lemma}\label{g2:lem:uniformly-supported-evolution}
	设 \(N\) 为无边界光滑流形，\(X_t\) 是 \(0\le t\le1\) 上光滑的时间依赖向量场。若存在紧集 \(L\subset N\)，使 \(X_t=0\) 于 \(N\setminus L\)，对所有 \(t\) 同时成立，则其演化映射 \(\Phi_{t,s}:N\to N\) 对全部 \(s,t\in[0,1]\) 有定义，光滑且为微分同胚，逆为 \(\Phi_{s,t}\)。它在 \(N\setminus L\) 上为恒等。
\end{lemma}
\begin{proof}
	第一卷第13章第1节“演化映射”给出局部存在、联合光滑性与逆映射公式。这里只需保证解能走到指定的时间。\(N\setminus L\) 中的点对全部时间都是零点，所以从这些点出发的解恒定。若一条从 \(L\) 出发的解进入 \(N\setminus L\)，从进入处反向使用唯一性，它在此前也应是同一条常值解，矛盾。因此这条解始终留在 \(L\)。

	若最大存在区间有一个位于 \([0,1]\) 内的有限端点，时空轨道就留在紧集 \([0,1]\times L\) 中。把 \(X_t\) 写成增广向量场 \(\partial_t+X_t\)，对这个紧集用有限个局部初值邻域取得统一的正存在时间，便能从靠近端点的时刻重新启动并越过它，与最大性矛盾。这正是第一卷第13章第2节逃逸判据的证明；端点0、1处先将光滑系数作局部时间延拓即可。故演化在整个指定时间区间存在，原有的逆映射与光滑依赖结论同时成立。
\end{proof}

\begin{theorem}[同位延拓定理]\label{g2:thm:isotopy-extension}
	设 \(K\) 为紧无边界光滑流形，\(N\) 为无边界光滑流形，\(i:[0,1]\times K\to N\) 为光滑同位。对任意包含 \(i([0,1]\times K)\) 的开集 \(\Omega\subset N\)，存在环境同位 \(\Phi_t\)，使
	\[
		\Phi_t\circ i_0=i_t,
	\]
	且所有 \(\Phi_t\) 在同一个紧集 \(L\Subset\Omega\) 外均为恒等。
\end{theorem}
\begin{proof}
	用上一小节的管状延拓与截断构造 \(X_t\)，再用\cref{g2:lem:uniformly-supported-evolution}取 \(\Phi_t=\Phi_{t,0}\)。对任意 \(x\in K\)，曲线 \(t\mapsto i_t(x)\) 满足
	\[
		\frac{\mathrm{d}}{\mathrm{d}t}i_t(x)=X_t(i_t(x)),\qquad i_{t=0}(x)=i_0(x).
	\]
	它与 \(t\mapsto\Phi_t(i_0(x))\) 是同一个初值问题的解，故由唯一性相等。共同紧支集给出所述支集控制。
\end{proof}

\begin{example}\label{g2:exm:translated-sphere-isotopy}
	设 \(a\in\mathbb R^{n+1}\)，\(i_t(x)=x+ta\)，其中 \(x\in S^n\)。整个运动轨迹紧，沿它的速度恒为 \(a\)。选择在轨迹邻域为1的紧支函数 \(\chi\)，以 \(X(y)=\chi(y)a\) 为环境向量场。每条指定轨道始终在 \(\chi=1\) 的区域，故环境流在球面上恰好实现 \(x\mapsto x+ta\)，在远处保持恒等。环境的其他点如何移动由截断决定，因此延拓存在而不唯一。
\end{example}

同位由嵌入组成这个要求不能放松成一般同伦。例如圆周上的常值映射与某些度为零的映射可以连续变形，但常值映射无法由环境微分同胚作用于一个嵌入得到；环境运动始终保留嵌入的单射性与局部秩。

\subsection{相对同位和非紧限制}
\label{g2:subsec:1505:04}

先看需要固定一个环境闭集 \(B\subset N\) 的情形。若存在开集 \(V\supset B\)，使轨迹进入 \(V\) 时的指定速度全部为零，即
\[
	\partial_ti(t,x)=0\qquad\text{当 }i_t(x)\in V,
\]
那么可让延拓后的环境同位在 \(B\) 的某邻域恒为恒等。事实上，速度非零的轨迹点的闭包是轨迹中的紧集，且避开 \(V\)；取一个在该紧集附近为1、支集避开 \(B\) 的截断函数，将前面的延拓场再乘以它。它在速度非零处保持原值，在其余轨迹点仍为零，所以仍延拓原同位，而支集与 \(B\) 分离。若仅要求同位在 \(K\) 的某个闭子集 \(A\) 上逐点不变，则原构造已经保证 \(X_t(i_0(x))=0\) 对 \(x\in A\) 成立，环境流因而逐点固定 \(i_0(A)\)；这并不自动固定其整个邻域。

\ProseParagraphBreak

紧源还可以换成受控的非紧源。一个足够条件是：\(i_0:K\to N\) 为固有嵌入，并存在固定紧集 \(C\subset K\)，使 \(i_t(x)=i_0(x)\) 对所有 \(x\notin C\) 和所有 \(t\) 成立。在这种情况下，\(\widehat i\) 是固有嵌入：紧目标集的原像落在 \([0,1]\times C\) 与 \([0,1]\times i_0^{-1}(L)\) 的并中，其中 \(L\) 为该目标集的空间投影。速度截面在轨迹上的支集紧，可在其附近按同样的局部标架构造延拓，并在其余轨迹上保留零值，再截断为有共同紧支集的场。流的论证不变，得到环境同位。

\begin{example}\label{g2:exm:noncompact-isotopy-no-extension}
	令 \(K=N=\mathbb R\)，定义
	\[
		i_t(x)=\begin{cases}
			\dfrac{\arctan(tx)}{t},&t>0,\\
			x,&t=0.
		\end{cases}
	\]
	积分表达 \(i_t(x)=\int_0^x(1+t^2u^2)^{-1}\,\mathrm{d}u\) 说明它在 \((t,x)\) 上光滑，且 \(\partial_x i_t>0\)，故每个 \(i_t\) 都是嵌入。当 \(t>0\) 时，其像是有界开区间 \((-\pi/(2t),\pi/(2t))\)；而 \(i_0\) 是整个实线的恒等映射。如果存在环境同位延拓它，\(\Phi_t\circ i_0=i_t\) 就迫使 \(\Phi_t=i_t\)，与 \(\Phi_t\) 必须满射到 \(\mathbb R\) 矛盾。非紧源上，各个切片是嵌入本身并不足够。
\end{example}

源 \(K\) 若紧但有边界，可以把速度截面的延拓改在半空间坐标中进行：先将各局部光滑分量延过边界，再用时空单位分解拼接，保留沿轨迹的原值。局部嵌入的半空间标准形由其延拓的浸入标准形给出，因此这一步适用于带边界的轨迹；截断和流的证明随后相同。环境 \(N\) 若有边界，所需延拓场必须沿 \(\partial N\) 切向，才能产生保持 \(N\) 的双侧运动。一个直接适用的相对版本是原同位在环境边界的某邻域内恒定，并把延拓场的支集选在内部；第一卷第13章第2节的边界切向流结论也说明了这一条件的作用。

% !TeX root = ../../Book2.tex
% 纲要标题：### 15.6 固有淹没的局部平凡性
% 纲要位置：几何学计划纲要.md，第 1610 行。
% 正文状态：初稿；尚未正式审核。

% 纲要内容开始
% 
% * 水平分布和局部提升方程
% * 固有性保证的统一存在区间
% * 径向提升与局部平凡化
% * Ehresmann 定理
% * Ehresmann 定理的带边界版本
% * 光滑与全纯局部平凡性
% 
% ---
% 
% 纲要内容结束

\section{固有淹没的局部平凡性}
\label{g2:sec:1506}

淹没定理说明，满秩映射在每个源点附近都像一个坐标投影。然而，一条纤维可能穿过许多张坐标图：把这些局部投影拼在一起，能否得到整个纤维附近的乘积？本节用沿基底路径移动纤维的方法回答这个问题。固有性保证移动过程中没有点在有限时间内逃向无穷远，因而同一套构造能作用于整条纤维。

\subsection{水平分布和局部提升方程}
\label{g2:subsec:1506:01}

先设 \(p:E\to B\) 是无边界光滑流形之间的淹没。记 \(\dim E=m\)、\(\dim B=b\)。每一点的核
\[
	V_e=\ker\mathrm{d}p_e\subseteq\mathrm{T}_eE
\]
有固定维数 \(m-b\)，由满秩矩阵的局部消元可知，它们组成 \(\mathrm{T}E\) 的光滑子丛 \(V\)，称为\term{垂直切丛}[vertical tangent bundle]。若 \(e\in p^{-1}(y)\)，则 \(V_e=\mathrm{T}_e(p^{-1}(y))\)：垂直方向就是留在当前纤维内的运动方向。

在 \(E\) 上选一个辅助 Riemann 度量，令 \(H=V^\perp\)。第一卷第12章已经证明，光滑向量丛可取光滑度量，光滑子丛的正交补仍是光滑子丛。因此
\[
	\mathrm{T}E=V\oplus H,
	\qquad \mathrm{d}p_e|_{H_e}:H_e\longrightarrow\mathrm{T}_{p(e)}B
\]
是线性同构。它的逆随 \(e\) 光滑变化：在局部标架中，这就是一个可逆光滑矩阵的逆。记这个逆为
\begin{equation}\label{g2:eq:horizontal-right-inverse}
	L_e:\mathrm{T}_{p(e)}B\longrightarrow\mathrm{T}_eE,
	\qquad \mathrm{d}p_e\circ L_e=\operatorname{id}.
\end{equation}
像丛 \(H\) 称为一个\term{水平分布}[horizontal distribution]。给定基底的一个速度，\(L_e\) 选出总空间中与它对应的速度。这里不要求水平分布满足 Frobenius 可积条件；提升一条曲线只需求解常微分方程，不需要存在处处切于 \(H\) 的高维子流形。

\ProseParagraphBreak

设 \(\psi:U_0\to\mathbb R^b\) 是基底坐标图，\(u\in\mathbb R^b\)。在 \(p^{-1}(U_0)\) 上定义依赖参数 \(u\) 的向量场
\begin{equation}\label{g2:eq:radial-horizontal-field}
	X_u(e)=L_e\bigl(\mathrm{d}(\psi^{-1})_{\psi(p(e))}u\bigr).
\end{equation}
若 \(e(t)\) 是它的积分曲线，链式法则给出
\[
	\frac{\mathrm{d}}{\mathrm{d}t}\psi(p(e(t)))
	=\mathrm{d}\psi_{p(e(t))}\mathrm{d}p_{e(t)}X_u(e(t))=u.
\]
所以基底坐标沿直线运动：
\begin{equation}\label{g2:eq:horizontal-base-line}
	\psi(p(e(t)))=\psi(p(e(0)))+tu.
\end{equation}
这条等式同时说明提升方程的含义与求解方向：我们先在基底指定直线，再用 \(L\) 确定总空间中的运动。

对一条光滑路径 \(c(t)\subset U_0\)，将上述常向量 \(u\) 换成 \(\frac{\mathrm{d}}{\mathrm{d}t}(\psi\circ c)(t)\)，便得到含时向量场。如果初始点位于 \(p^{-1}(c(0))\)，同一个计算说明解始终投影到 \(c(t)\)。第一卷第13章第1节的局部存在唯一性和初值、参数光滑依赖定理适用：系数关于 \((t,e)\) 光滑，而含时方程也可通过增加时间坐标写成自治方程。下面需要补上的，是这些局部解能存在多久。

\subsection{固有性保证的统一存在区间}
\label{g2:subsec:1506:02}

回忆，\(p\) 固有是指每个紧集 \(K\subset B\) 的原像 \(p^{-1}(K)\) 紧。因而，只要基底路径始终留在同一个紧集内，它的任何水平提升也都被限制在一个紧集内。这与第一卷第13章第2节的逃逸判据正好配合：光滑向量场的最大积分曲线若在有限时刻不能继续，就必定离开定义域中的每个紧集。

\begin{lemma}\label{g2:lem:proper-horizontal-lift}
	设 \(p:E\to B\) 是固有淹没，\(c:[0,1]\to U_0\) 是光滑路径，且 \(c([0,1])\) 包含在 \(U_0\) 的某个紧子集中。对每个 \(e_0\in p^{-1}(c(0))\)，由选定水平分布确定的提升都在整个 \([0,1]\) 上存在且唯一。
\end{lemma}
\begin{proof}
	先由局部存在唯一性得到经过 \(e_0\) 的提升。设它的向前最大存在时刻为 \(T\leq1\)。在 \(t<T\) 时，它的投影等于 \(c(t)\)，所以提升位于紧集 \(p^{-1}(c([0,1]))\subset p^{-1}(U_0)\) 内。把时间作为附加坐标后，积分曲线位于紧集 \([0,1]\times p^{-1}(c([0,1]))\) 内。路径及向量场可在端点附近光滑延拓；有限个局部存在邻域给出这个紧集附近统一的短时存在区间。因此解可以越过 \(T\) 继续，与 \(T\) 是不可延拓的端点矛盾。若只需解到 \(T=1\)，同一论证保证端点值存在。唯一性来自局部唯一性。
\end{proof}

这里的“统一”不要求预先给整条非紧纤维估计一个全局速度界。固有性已经使初始纤维和路径上方的全部区域紧，从而可用有限个局部方程控制所有初始点。对于依赖光滑参数的路径，还可把参数加入自治方程，令其导数为零；当参数限制在紧集内时，同样的有限覆盖论证给出共同的存在区间。随后，初值与参数光滑依赖定理便保证提升同时光滑依赖于起点和参数。关于这两个常微分方程工具，可对照 \cite[第9章，第215--217页]{Lee2012Smooth}。

\ProseParagraphBreak

固有性在这里有实际作用。例如
\[
	p:\mathbb R^2\setminus\{(0,0)\}\longrightarrow\mathbb R,
	\qquad p(x,y)=x
\]
处处为淹没，但 \(p^{-1}(0)\cong\mathbb R\setminus\{0\}\) 有两个连通分支，\(p^{-1}(x)\cong\mathbb R\) 对 \(x\ne0\) 只有一个连通分支。因此 \(0\) 附近不存在以整条纤维为因子的局部乘积。满秩条件没有阻止纤维的一部分从缺失点附近消失；这个映射也确实不固有。

\subsection{径向提升与局部平凡化}
\label{g2:subsec:1506:03}

固定 \(y_0\in B\)，取坐标图 \(\psi:U_0\to\mathbb R^b\)，使 \(\psi(y_0)=0\)。缩小后可取 \(r>0\)，使闭球 \(\overline{B(0,2r)}\) 包含在 \(\psi(U_0)\) 中。令
\[
	U=\psi^{-1}(B(0,r)),\qquad F=p^{-1}(y_0).
\]
对 \(y\in U\)，选择从 \(y_0\) 到 \(y\) 的坐标径向路径
\[
	c_y(t)=\psi^{-1}(t\psi(y)),\qquad 0\leq t\leq1.
\]
每一条路径都位于紧集 \(K=\psi^{-1}(\overline{B(0,r)})\) 内，而且随 \((t,y)\) 光滑变化。

记 \(\varphi_u^t(e)\) 为式 \eqref{g2:eq:radial-horizontal-field} 中向量场 \(X_u\) 的局部流。由 \cref{g2:lem:proper-horizontal-lift}，对任意 \((z,y)\in F\times U\)，表达式
\begin{equation}\label{g2:eq:ehresmann-trivialization}
	\Phi(z,y)=\varphi_{\psi(y)}^1(z)
\end{equation}
都有定义。式 \eqref{g2:eq:horizontal-base-line} 说明
\[
	p(\Phi(z,y))=y.
\]
因此 \(\Phi\) 已经具有局部平凡化所需的投影关系。它把纤维 \(F\) 中每一点沿同一条基底路径搬运到 \(p^{-1}(y)\)，而不是分别在源空间中挑选互不相干的坐标图。

为了从另一条纤维回到 \(F\)，对 \(e\in p^{-1}(U)\) 固定 \(u=\psi(p(e))\)，沿 \(X_u\) 反向流动一个单位时间。它在基底中的坐标是 \((1+t)u\)，\(-1\leq t\leq0\)，仍留在 \(K\) 内，故同样存在。于是候选逆映射为
\begin{equation}\label{g2:eq:ehresmann-inverse}
	\Psi(e)=\bigl(\varphi_{\psi(p(e))}^{-1}(e),p(e)\bigr).
\end{equation}
这里上标 \(-1\) 表示时间 \(-1\)，而参数 \(\psi(p(e))\) 在这次反向流动中保持不变。若在途中不断按当前位置重新设定参数，方程就变了，也不再是原路径的逆向提升。

\subsection{Ehresmann 定理}
\label{g2:subsec:1506:04}

\begin{theorem}[Ehresmann 定理]\label{g2:thm:ehresmann}
	设 \(E,B\) 是无边界光滑流形，\(p:E\to B\) 是固有的满淹没。则每个 \(y_0\in B\) 都有邻域 \(U\)，使得存在微分同胚
	\[
		\Phi:F\times U\longrightarrow p^{-1}(U),\qquad F=p^{-1}(y_0),
	\]
	满足 \(p\circ\Phi=\operatorname{pr}_U\)。特别地，\(p\) 是以紧流形为纤维的光滑纤维丛。
\end{theorem}
\begin{proof}
	淹没定理给出 \(F\) 的光滑子流形结构，固有性给出其紧性。按照前两小节选取水平分布和坐标径向路径，并定义式 \eqref{g2:eq:ehresmann-trivialization} 与 \eqref{g2:eq:ehresmann-inverse}。

	先验证光滑性。向量场 \(X_u(e)\) 关于 \((e,u)\) 光滑，所以流映射 \((t,e,u)\mapsto\varphi_u^t(e)\) 在其存在域内光滑。它的存在域是开集，这是初值与参数光滑依赖定理的一部分。上一小节已证明，定义 \(\Phi\) 所需的每个 \((1,z,\psi(y))\) 和定义 \(\Psi\) 所需的每个 \((-1,e,\psi(p(e)))\) 都在这个存在域内。因此限制到这些初值并复合光滑映射后，\(\Phi\) 与 \(\Psi\) 都光滑。此处分别使用了两个事实：固有性保证解确实走到指定时刻，参数依赖定理保证到达点光滑变化。

	再验证互逆。若 \(u=\psi(y)\)，则 \(p(\varphi_u^1(z))=y\)，故逆向运动使用的仍是同一个向量场 \(X_u\)。流的唯一性给出
	\[
		\varphi_u^{-1}(\varphi_u^1(z))=z,
		\qquad \varphi_u^1(\varphi_u^{-1}(e))=e.
	\]
	于是 \(\Psi\Phi=\operatorname{id}_{F\times U}\)、\(\Phi\Psi=\operatorname{id}_{p^{-1}(U)}\)。投影关系已经验证，定理得证。
\end{proof}

若基底连通，则它道路连通。沿连接两个基点的路径取有限个局部平凡化，逐次比较相邻纤维，就知所有纤维彼此微分同胚。这个结论仍允许全局的粘合发生扭转。例如 \(S^1\to S^1\)、\(z\mapsto z^2\) 是固有满淹没，纤维为两个点，局部有上述乘积；总空间 \(S^1\) 连通，而 \(\{1,2\}\times S^1\) 不连通，因此全局并不是这个乘积。

所选度量、水平分布和基底径向坐标都影响具体的 \(\Phi\)。更换选择后仍得到合法的局部平凡化；两种平凡化之间的坐标变换在每条纤维上是一个微分同胚，并随基点光滑变化。定理保证局部乘积的存在，不指定一个天然的纤维识别。

\subsection{Ehresmann 定理的带边界版本}
\label{g2:subsec:1506:05}

有边界时，水平提升还必须尊重边界。先讨论基底无边界、总空间可能有边界的情形。

\begin{proposition}\label{g2:prop:ehresmann-source-boundary}
	设 \(E\) 是有边界光滑流形，\(B\) 无边界，\(p:E\to B\) 固有、满射且处处为淹没。若限制 \(p|_{\partial E}:\partial E\to B\) 也是淹没，则每个 \(y_0\in B\) 附近都有保持边界的局部平凡化
	\[
		(p^{-1}(U),p^{-1}(U)\cap\partial E)
		\cong(F,\partial F)\times U,
		\qquad \partial F=F\cap\partial E.
	\]
\end{proposition}
\begin{proof}
	关键是构造满足 \eqref{g2:eq:horizontal-right-inverse} 的 \(L\)，并使边界点的 \(L_e\) 取值于 \(\mathrm{T}_e\partial E\)。在边界点附近，\(p|_{\partial E}\) 满秩，故可将基底坐标 \(q^1,\ldots,q^b\) 的拉回补成边界坐标，再加一个边界定义函数 \(r\geq0\)。这些函数给出适配坐标 \((q,z,r)\)，其中 \(p\) 就是 \(q\) 投影。令基底坐标方向的局部提升保持 \((z,r)\) 不变，便得到沿边界切向的局部右逆。内点附近按普通淹没坐标选择局部右逆。

	取从属于这些邻域的光滑单位分解 \(\{\rho_i\}\)，把局部右逆 \(L_i\) 拼成 \(L=\sum_i\rho_iL_i\)。由于 \(\sum_i\rho_i=1\)，有 \(\mathrm{d}p\circ L=\operatorname{id}\)；在边界点，每个有贡献的 \(L_i\) 都取值于边界切空间，因此 \(L\) 也如此。其像与 \(\ker\mathrm{d}p\) 互补。

	现在按前述公式构造 \(X_u\)。它沿 \(\partial E\) 切向，其流保持边界，也保持内部。第一卷第13章第2节已说明这一点：在半空间坐标中，切向条件使法向方程写成 \(\dot r=a(q,z,r)r\)，所以 \(r=0\) 与 \(r>0\) 都在存在区间内保持。紧集上的统一短时存在和逃逸判据也适用于这样的向量场。故固有性保证正、反向提升均存在。原来的 \(\Phi\) 与 \(\Psi\) 保持边界且互逆，给出结论。适配坐标同时说明纤维是有边界子流形，边界恰为 \(F\cap\partial E\)。
\end{proof}

不能直接对任意度量取 \((\ker\mathrm{d}p)^\perp\)，就宣称所得水平分布沿边界切向：正交性并不自动带来这个条件。上面的局部右逆加单位分解，正是为解决这一点而作。

\ProseParagraphBreak

边界上的满秩条件也不能删除。考虑光滑带边界曲面
\[
	E=\{(x,y)\in\mathbb R^2:|y|\leq2,\ x^2+y^2\geq1\},
	\qquad p(x,y)=x.
\]
它的边界由两条水平直线和单位圆组成，彼此不相交，因而没有角点。映射 \(p\) 固有、满射，而且 \(\mathrm{d}p=\mathrm{d}x\) 在二维切空间上处处满射。然而，\(|x|<1\) 时纤维为两段闭区间，\(|x|>1\) 时纤维为一段闭区间，所以在 \(x=\pm1\) 附近不可能局部平凡。问题恰发生在圆周上的 \((\pm1,0)\)：此处 \(p|_{\partial E}\) 的微分为零。

\ProseParagraphBreak

如果基底也有边界，需说明总空间的哪一部分对应基底边界。一个适用于无角点流形的清楚版本如下。

\begin{proposition}\label{g2:prop:ehresmann-base-boundary}
	设 \(E,B\) 是有边界光滑流形，\(p:E\to B\) 是固有满淹没，并满足
	\[
		p^{-1}(\partial B)=\partial E.
	\]
	则 \(p\) 在基底的半空间邻域上也有光滑局部平凡化，纤维是紧的无边界流形。
\end{proposition}
\begin{proof}
	先看 \(e\in\partial E\)。取源和目标的边界坐标，分别记法向坐标为 \(r\geq0\)、\(s\geq0\)。因为 \(s\circ p\) 在 \(r=0\) 上为零，可写成
	\[
		s(p(q,r))=r\,a(q,r).
	\]
	其边界切向导数为零。满秩条件迫使 \(a(q,0)\ne0\)，而内点映到内点迫使 \(a(q,0)>0\)。因此 \(p\) 在法向上为同向的非退化映射；同时，\(\mathrm{d}p\) 将边界切空间满射到基底的边界切空间。于是 \(p|_{\partial E}:\partial E\to\partial B\) 是固有满淹没。

	固定 \(y_0\in\partial B\)。对这个无边界流形之间的限制映射应用 \cref{g2:thm:ehresmann}，得到某个边界邻域 \(V\) 上的平凡化，纤维 \(F=p^{-1}(y_0)\) 紧且无边界。把 \(V\) 缩小，使其闭包位于一张边界坐标图中，再取基底的半空间邻域 \(V\times[0,\varepsilon)\)，记其法向坐标为 \(s\)。在该邻域上选择 \(\mathrm{d}p\) 的光滑右逆，提升 \(\partial/\partial s\) 得到向量场 \(Y\)。上面的法向计算说明，\(Y\) 在 \(\partial E\) 上严格指向内侧。

	从边界平凡化中的点出发，沿 \(Y\) 流动时间 \(s\)。基底的 \(V\) 坐标保持不变，法向坐标恰好由 \(0\) 增为 \(s\)。紧性给出统一的短时存在，缩小 \(\varepsilon\) 后定义出
	\[
		F\times V\times[0,\varepsilon)\longrightarrow p^{-1}(V\times[0,\varepsilon)).
	\]
	任意右侧点的逆构造是沿 \(Y\) 反向流动，直到法向坐标降为零，再使用边界平凡化的逆。在到达零之前，轨道保持在内部；其基底轨道是紧线段，固有性防止在途中逃逸。到达边界时，向量场在局部半空间坐标中可光滑延拓到开集，故流与反向到达映射一直光滑到边界。唯一性证明这两个构造互逆。基底内点附近则直接应用无边界版本。
\end{proof}

例如闭圆盘上的投影 \(\overline D^2\to[-1,1]\)、\((x,y)\mapsto x\) 固有、满射，在二维切空间上处处满秩，但它不满足 \(p^{-1}(\partial B)=\partial E\)。端点上方的纤维为单点，内部纤维为线段，确实不能局部平凡。若还希望纤维本身有边界，那么它与有边界基底的乘积通常会出现角点；此时需另行规定各边界面的对应关系，上述无角点版本不涉及这种更一般的结构。

\subsection{光滑与全纯局部平凡性}
\label{g2:subsec:1506:06}

一个复流形同时是实光滑流形。全纯映射的坐标表达满足复可微条件，而一般的实光滑映射没有这个约束。因此，固有全纯淹没可以应用 \cref{g2:thm:ehresmann}，但所得平凡化首先是实光滑的。下面用复环面说明，全纯平凡化确实可能不存在。

取上半平面 \(\{\tau\in\mathbb C:\operatorname{Im}\tau>0\}\) 中的一个小开圆盘 \(U\)，定义格子与复环面
\[
	\Lambda_\tau=\mathbb Z+\mathbb Z\tau,
	\qquad E_\tau=\mathbb C/\Lambda_\tau.
\]
把这些环面组成一个族，可令 \(\mathbb Z^2\) 作用于 \(\mathbb C\times U\)：
\[
	(m,n)\cdot(z,\tau)=(z+m+n\tau,\tau),
\]
并取商空间 \(E=(\mathbb C\times U)/\mathbb Z^2\)。这个作用自由且适当不连续。后一条可以直接检查：在任意紧集上，\(\operatorname{Im}\tau\) 有正下界，而 \(z\) 有界；若平移前后的点都在该紧集中，先由虚部估计得到整数 \(n\) 有界，再由实部估计得到 \(m\) 有界，故只有有限个群元素能使紧集与其像相交。由于作用由全纯变换给出，局部商坐标使 \(E\) 成为复流形，映射
\[
	p:E\longrightarrow U,\qquad [z,\tau]\longmapsto\tau
\]
是全纯淹没，纤维正是 \(E_\tau\)。

这个族甚至有一个显式的实光滑全局平凡化：
\begin{equation}\label{g2:eq:complex-torus-smooth-family}
	(\mathbb R^2/\mathbb Z^2)\times U\longrightarrow E,
	\qquad ([s,t],\tau)\longmapsto[s+t\tau,\tau].
\end{equation}
它的逆在代表元上由
\[
	t=\frac{\operatorname{Im}z}{\operatorname{Im}\tau},
	\qquad s=\operatorname{Re}z-t\operatorname{Re}\tau
\]
给出；更换格子代表元只改变 \((s,t)\) 的整数部分。因此它是良定义的实光滑微分同胚。纤维 \(\mathbb R^2/\mathbb Z^2\) 紧，故 \(p\) 固有。

\ProseParagraphBreak

这些实光滑同胚的纤维是否也彼此双全纯？这里“双全纯”是指映射及其逆都全纯。下面的判别将把这个问题归结为格子的变化。

\begin{lemma}\label{g2:lem:complex-torus-lattices}
	若 \(\mathbb C/\Lambda\) 与 \(\mathbb C/\Lambda'\) 双全纯，则存在 \(a\in\mathbb C\setminus\{0\}\)，使 \(a\Lambda=\Lambda'\)。反过来，这个条件也给出双全纯同构。
\end{lemma}
\begin{proof}
	商映射 \(\mathbb C\to\mathbb C/\Lambda\) 是万有覆盖。给定双全纯同构 \(f\)，把它与前一个覆盖映射复合。由于 \(\mathbb C\) 单连通，覆盖映射的提升定理给出提升 \(F:\mathbb C\to\mathbb C\)；覆盖映射局部双全纯，所以 \(F\) 也是全纯的。这里用到 \cref{g2:prop:covering-lift-criterion}：\(\mathbb C\) 道路连通、局部道路连通，且基本群为零，因此满足提升所需的子群条件。

	对每个 \(\lambda\in\Lambda\)，差 \(F(z+\lambda)-F(z)\) 取值于离散集 \(\Lambda'\)，又随 \(z\) 连续，故为常量。求导得 \(F'(z+\lambda)=F'(z)\)。于是 \(F'\) 在一个紧基本平行四边形上有界，周期性又使它在全平面上有界。由有界整函数为常数的 Liouville 定理，\footnote{Joseph Liouville（1809--1882），法国数学家。}\(F'=a\)，从而 \(F(z)=az+c\)。这里使用的复分析结论可由 Cauchy 导数估计迅速推出。\footnote{Augustin-Louis Cauchy（1789--1857），法国数学家，研究极限与收敛的严格表述，建立了复函数积分理论的基本结果。}若整函数 \(h\) 满足 \(|h|\leq M\)，则在任意半径 \(R\) 的圆盘中心有 \(|h'|\leq M/R\)，令 \(R\to\infty\) 即得 \(h'=0\)。

	因为 \(f\) 为局部同构，\(a\ne0\)；又因 \(F(z+\lambda)-F(z)=a\lambda\in\Lambda'\)，有 \(a\Lambda\subseteq\Lambda'\)。若某个 \(\mu\in\Lambda'\) 不属于 \(a\Lambda\)，则 \(z\) 与 \(z+\mu/a\) 在 \(\mathbb C/\Lambda\) 中不同，却被 \(f\) 映到同一点，与单射性矛盾。故 \(a\Lambda=\Lambda'\)。反过来，\(z\mapsto az\) 在商空间上直接诱导所需同构。
\end{proof}

对 \(\Lambda_\tau=\mathbb Z+\mathbb Z\tau\)，上述格子条件等价于
\begin{equation}\label{g2:eq:torus-modular-equivalence}
	\sigma=\frac{A\tau+B}{C\tau+D},
	\qquad
	\begin{pmatrix}A&B\\C&D\end{pmatrix}\in\mathrm{SL}(2,\mathbb Z).
\end{equation}
理由是，\((1,\sigma)\) 与 \((a,a\tau)\) 是同一个格子的两组正定向整数基，换基矩阵的行列式为 \(1\)；将两个基的关系相除便得到此式。反过来，若有此式，则 \((C\tau+D,A\tau+B)\) 是 \(\Lambda_\tau\) 的整数基，乘以 \((C\tau+D)^{-1}\) 后得到 \(\Lambda_\sigma\)。

固定一个 \(\tau_0\)，整数矩阵只有可数多个，所以与 \(E_{\tau_0}\) 双全纯的参数至多可数。任意开圆盘却有不可数多个参数，故不存在 \(\tau_0\) 的邻域，使该族中的所有纤维都与 \(E_{\tau_0}\) 双全纯。全纯局部平凡化会给出这样的纤维同构，因此不可能存在。

式 \eqref{g2:eq:complex-torus-smooth-family} 将差别呈现得很具体：所有纤维的实光滑结构相同，格子的复形状却随 \(\tau\) 改变。Ehresmann 证明中的辅助度量、水平补空间和单位分解都是实光滑构造，它们正好提供前一种意义的平凡化。


\begin{chaptersummary}
	坐标与截断给出初始嵌入，切向和割线坏方向的维数估计给出弱 Whitney 界；非紧源还要保持固有性，才能由单射浸入得到嵌入。强界要求继续处理一般位置留下的双点，高维 Whitney 移动、低维模型和非紧源构造分别承担不同任务。法向局部指数在零截面处的微分可逆，变半径控制把这些局部逆拼成管状邻域；边界的向内向量场则给出领圈，为双化与粘合提供横向坐标。同位延拓把嵌入族的速度扩展到环境，再由共同紧支集保证环境流存在。Ehresmann 定理使用另一种存在性控制：固有性把紧基底路径的提升限制在紧集内，径向提升及其反向过程于是给出光滑局部平凡化。源有边界时还须控制边界限制的秩；即使已获光滑平凡化，复结构也可以随参数变化。
\end{chaptersummary}

\BookExercises

\begin{exercise}\label{g2:ex:projection-loses-properness}
	设 \(F:\mathbb R\to\mathbb R^2\)、\(F(t)=(\arctan t,t)\)，令 \(p:\mathbb R^2\to\mathbb R\) 为第一坐标投影。证明 \(F\) 是固有嵌入，而 \(p\circ F\) 是非固有的嵌入。将 \(F(\mathbb R)\) 包含在竖直轴周围一个有界半径管中，并说明本例中\cref{g2:lem:proper-projection-tube}的哪个假设不成立。
\end{exercise}
\begin{solution}
	\(F'(t)=(1/(1+t^2),1)\ne0\)，第二坐标使 \(F\) 单射，且它在像上的连续逆就是第二坐标限制，故为嵌入。若 \(K\subset\mathbb R^2\) 紧，则 \(F^{-1}(K)\) 闭，并被 \(K\) 的第二坐标界限制在某个有界区间内，故紧。因此 \(F\) 固有。

	\(p\circ F=\arctan\) 的导数恒正，它是 \(\mathbb R\) 到 \((-\pi/2,\pi/2)\) 的微分同胚，所以作为到 \(\mathbb R\) 的映射仍为嵌入。但紧区间 \([-\pi/2,\pi/2]\) 的原像为整个 \(\mathbb R\)，故不固有。由于 \(|\arctan t|<\pi/2\)，原曲线处于竖直轴周围半径 \(\pi/2\) 的管中；本次投影沿竖直方向，恰好压掉了管轴，因此引理所要求的投影方向不平行管轴的条件失效。
\end{solution}

\begin{exercise}\label{g2:ex:proper-immersion-finite-branches}
	设 \(f:M\to N\) 是无边界光滑流形间的固有浸入。证明每个纤维 \(f^{-1}(y)\) 有限。若 \(f^{-1}(y)=\{x_1,\ldots,x_s\}\)，证明存在 \(y\) 的开邻域 \(V\) 以及两两不交的 \(x_i\) 的开邻域 \(U_i\)，使 \(f^{-1}(V)\subset\bigcup_iU_i\)，且每个 \(f|_{U_i}\) 为嵌入。这说明像在 \(y\) 附近有多少条局部分支？
\end{exercise}
\begin{solution}
	常秩定理说明浸入在每个源点附近都是嵌入，故纤维是离散集。固有性又保证单点的原像紧。若紧离散纤维无限，其逐点孤立邻域构成没有有限子覆盖的开覆盖，矛盾，故纤维有限。

	在有限多个 \(x_i\) 附近取两两不交的开集 \(U_i\)，使各限制为嵌入。若不存在所需 \(V\)，可选 \(y_j\to y\) 及 \(z_j\notin\bigcup_iU_i\)，满足 \(f(z_j)=y_j\)。集合 \(\{y\}\cup\{y_j:j\ge1\}\) 紧，其固有原像也紧，因此 \(z_j\) 有子列收敛到 \(z\)。连续性给出 \(f(z)=y\)，故 \(z=x_i\) 对某个 \(i\) 成立；但 \(U_i\) 开，子列最终应落入 \(U_i\)，矛盾。于是所需 \(V\) 存在。像在 \(V\) 中由 \(s\) 条嵌入分支组成；不同分支可能相交，也可能局部重合，因此这不等于它们有 \(s\) 个不同的像。
\end{solution}

\begin{exercise}\label{g2:ex:projective-rank-model}
	设 \(m\ge2\)，\(V_m\) 为实对称无迹 \((m+1)\times(m+1)\) 矩阵空间，\(A_m=\operatorname{diag}(1,\ldots,1,-m)\)，\(P_m:V_m\to V_m/\mathbb RA_m\) 为商映射。证明
	\[
		[x]\longmapsto P_m\left(xx^{\mathsf T}-\frac1{m+1}I\right),\qquad |x|=1,
	\]
	给出 \(\mathbb RP^m\) 到维数为 \((m+1)(m+2)/2-2\) 的 Euclid 空间的光滑嵌入。判断哪些 \(m\ge2\) 使这一维数等于强 Whitney 界 \(2m\)，并解释同一个公式在 \(m=1\) 时为何失效。
\end{exercise}
\begin{solution}
	对称矩阵空间的维数是 \((m+1)(m+2)/2\)，无迹条件及取一维商各减一。公式在 \(x\) 和 \(-x\) 处相同，故在射影空间上良定义。若两点像相等，则 \(xx^{\mathsf T}-yy^{\mathsf T}=\lambda A_m\)。左端秩至多 \(2\)，右端在 \(\lambda\ne0\) 时秩为 \(m+1\ge3\)，故 \(\lambda=0\)。于是两秩一正交投影相同，像直线相同，得到 \([x]=[y]\)。

	对 \(v\perp x\)，微分的代表为 \(xv^{\mathsf T}+vx^{\mathsf T}\)。它仍有秩至多 \(2\)，故若在商中为零，就必须为零；把此矩阵作用于 \(x\)，得到 \(v=0\)。所以映射是单射浸入。射影空间紧，因而它是嵌入。

	等式 \((m+1)(m+2)/2-2=2m\) 化为 \((m-2)(m+1)=0\)，在所给范围只有 \(m=2\)。当 \(m=1\) 时，取 \(x=(1,0)\)、\(y=(0,1)\)，便有 \(xx^{\mathsf T}-yy^{\mathsf T}=A_1\)，两个不同射影点在商中重合；秩比较失效的原因是 \(A_1\) 本身只有秩 \(2\)。
\end{solution}

\begin{exercise}\label{g2:ex:projection-many-double-points}
	对正整数 \(r\)，设
	\[
		\gamma_r(t)=(\sin t,\sin(2rt),\cos t),\qquad t\in\mathbb R/2\pi\mathbb Z.
	\]
	证明 \(\gamma_r\) 是光滑嵌入，而其前两个坐标组成的平面投影是恰有 \(2r-1\) 个横截双点的浸入，且没有三重点。
\end{exercise}
\begin{solution}
	第一、第三坐标组成标准圆周参数，所以 \(\gamma_r\) 单射且微分不为零；源紧，故是嵌入。投影的微分为 \((\cos t,2r\cos(2rt))\)。若 \(\cos t=0\)，则 \(t=\pi/2\) 或 \(3\pi/2\) 模 \(2\pi\)，此时 \(\cos(2rt)=(-1)^r\ne0\)，故投影始终为浸入。

	若两个不同参数 \(t,u\) 的投影相同，第一坐标给出 \(u=\pi-t\) 模 \(2\pi\)，而第二坐标变成 \(\sin(2rt)=-\sin(2rt)\)。所以 \(t=k\pi/(2r)\)，其中 \(0\le k<4r\)。排除对合 \(t\mapsto\pi-t\) 的两个不动点 \(\pi/2,3\pi/2\)，其余 \(4r-2\) 个参数两两成对，故共有 \(2r-1\) 个双点。一个给定的正弦值在圆周上至多有两个原像，所以不会出现三重点。

	在配对的 \(t,u=\pi-t\) 处，两条切向量分别为
	\[
		(\cos t,2r\cos(2rt)),\qquad(-\cos t,2r\cos(2rt)).
	\]
	其行列式为 \(4r\cos t\cos(2rt)\)。候选参数处 \(\cos(2rt)=\pm1\)，而 \(\cos t=0\) 的两个参数已经排除，故所有双点都横截。
\end{solution}

\begin{exercise}\label{g2:ex:normal-trivial-defining-map}
	设 \(S\subset N\) 是余维 \(c\) 的无边界嵌入子流形。证明：法丛 \(\nu_NS\) 平凡，当且仅当存在 \(S\) 的开邻域 \(U\subset N\) 和光滑淹没 \(q:U\to\mathbb R^c\)，使 \(q^{-1}(0)=S\)。这里的方程只要求在一个邻域内成立。
\end{exercise}
\begin{solution}
	若存在这样的 \(q\)，正则水平集的切空间公式给出 \(\ker\mathrm{d}q_x=\mathrm{T}_xS\)。因此 \(\mathrm{d}q_x\) 在法商上诱导线性同构 \(\nu_xS\to\mathbb R^c\)。该同构光滑依赖 \(x\)，于是 \([v]\mapsto(x,\mathrm{d}q_xv)\) 是 \(\nu_NS\cong S\times\mathbb R^c\) 的丛平凡化。

	反过来，取平凡化 \(A:\nu_NS\to S\times\mathbb R^c\) 和管状映射 \(\tau:E_\varepsilon\to U\)。令
	\[
		q=\operatorname{pr}_2\circ A\circ\tau^{-1}:U\to\mathbb R^c.
	\]
	\(A(E_\varepsilon)\) 是乘积流形中的开集，而第二坐标投影在该开集上的微分满射，所以 \(q\) 是淹没。它的零点对应法丛零截面，正好是 \(S\)。证明同时说明：一个一般子流形不能用同一组独立方程沿全体基点描述，障碍可能来自法丛的拼接。
\end{solution}

\begin{exercise}\label{g2:ex:diagonal-tubular-model}
	设 \(M\) 为无边界光滑流形，\(\Delta_M\subset M\times M\) 为对角线。用 \([(u,v)]\mapsto v-u\) 证明 \(\nu_{M\times M}\Delta_M\cong\mathrm{T}M\)。选择 \(\mathrm{T}M\) 上一个联络，证明
	\[
		\Psi(x,w)=(x,\exp_x^\nabla w)
	\]
	在 \(\mathrm{T}M\) 的零截面某个变半径邻域上给出对角线的管状邻域映射。交换两个因子的映射在对角线的法丛上诱导什么作用？
\end{exercise}
\begin{solution}
	在 \((x,x)\) 处，\(\mathrm{T}_{(x,x)}(M\times M)=\mathrm{T}_xM\oplus\mathrm{T}_xM\)，而对角线的切空间为 \(\{(u,u)\}\)。线性映射 \((u,v)\mapsto v-u\) 满射，核恰为对角切空间，因此诱导法商同构；它由切向量的自然运算定义，随 \(x\) 光滑，给出所需丛同构。

	\(\Psi(x,0)=(x,x)\)。在零截面处，由局部指数的微分公式得
	\[
		\mathrm{d}\Psi_{(x,0)}(u,w)=(u,u+w),
	\]
	这是线性同构，且其法商分量为 \((u+w)-u=w\)。对角线在 Hausdorff 乘积中闭，故\cref{g2:lem:zero-section-shrinking}直接给出一个变半径邻域，使 \(\Psi\) 成为到开像的微分同胚。以上法商计算又保证正规化条件，所以它是管状映射。

	交换映射把 \((u,v)\) 变成 \((v,u)\)，法商中的差由 \(v-u\) 变为 \(u-v\)，因此所诱导的作用是 \(\mathrm{T}M\) 各纤维上的负恒等。
\end{solution}

\begin{exercise}\label{g2:ex:tubular-normalization-sign}
	对 \(S=\{0\}\subset\mathbb R\)，考虑零点附近的两张邻域参数化 \(\tau_+(v)=v\) 和 \(\tau_-(v)=-v\)。证明不存在共同零点邻域上的联合光滑族 \(\tau_t\)，使各 \(\tau_t\) 都是局部微分同胚、固定零点，且端点分别为这两张参数化。解释为何这不违反\cref{g2:prop:tubular-uniqueness}，并求缩放公式 \(\delta_{1/t}\tau_-\delta_t\) 在 \(t\to0\) 时的极限。
\end{exercise}
\begin{solution}
	若这样的族存在，函数 \(a(t)=\mathrm{d}(\tau_t)_0\) 连续且处处非零，因为每张映射在零点都是局部微分同胚。然而 \(a(0)=1\)、\(a(1)=-1\)，介值定理迫使某个 \(a(t)=0\)，矛盾。

	这里法丛就是 \(\mathbb R\)，\(\tau_-\) 在法商上诱导 \(-\operatorname{id}\)，没有满足管状映射定义中的恒等正规化条件，故唯一性命题不适用。对 \(t>0\)，直接计算
	\[
		\delta_{1/t}\circ\tau_-\circ\delta_t(v)=\frac{-tv}{t}=-v,
	\]
	极限仍为负恒等，并非恒等。若先把纤维识别乘以 \(-1\)，则参数化变成 \(v\mapsto v\)，这时两者才具有同一正规化。
\end{solution}

\begin{exercise}\label{g2:ex:boundary-functions-ratio}
	设 \(r_0,r_1:M\to[0,\infty)\) 都是光滑带边界流形 \(M\) 的边界定义函数，即零点集均为 \(\partial M\)，且边界上的微分均非零。证明在边界的某个开邻域中存在正光滑函数 \(a\)，使 \(r_1=ar_0\)，并求二者在边界上的微分关系。
\end{exercise}
\begin{solution}
	按\cref{g2:lem:boundary-defining-function}的归一化步骤，在边界附近取 \(Xr_0=1\) 的向内向量场。领圈证明中的流参数化 \(F(x,t)\) 满足 \(r_0(F(x,t))=t\)，其中 \(t\) 可在一个随 \(x\) 变化的正区间内取值。令 \(b(x,t)=r_1(F(x,t))\)。因为 \(b(x,0)=0\)，微积分基本定理给出
	\[
		b(x,t)=t\int_0^1\partial_tb(x,ut)\,\mathrm{d}u.
	\]
	积分定义了包括 \(t=0\) 在内的光滑函数。\(r_1\) 在内部为正且边界微分非零，所以 \(\partial_tb(x,0)>0\)；缩小邻域后，该积分仍正。通过 \(F\) 的逆把它搬回 \(M\)，便得到所需 \(a\)。对 \(r_1=ar_0\) 求微分，并用边界上的 \(r_0=0\)，得到 \(\mathrm{d}r_1=a\,\mathrm{d}r_0\)。正因这个比例为正，由任一边界定义函数判断的向内方向都相同。
\end{solution}

\begin{exercise}\label{g2:ex:interior-homotopy-equivalence}
	设 \(M\) 为光滑带边界流形，\(j:\operatorname{Int}M\hookrightarrow M\) 为包含映射。用领圈证明 \(j\) 是同伦等价；允许 \(M\) 及其边界非紧。
\end{exercise}
\begin{solution}
	取领圈 \(c:\partial M\times[0,1)\to M\)。选择非负光滑函数 \(\chi:[0,1)\to[0,1]\)，使它在零附近为1，在 \([1/2,1)\) 上为零。再选 \(\varepsilon>0\)，使 \(\varepsilon\sup|\chi'|<1/2\)。对 \(0\le s\le1\) 定义
	\[
		f_s(t)=t+s\varepsilon\chi(t).
	\]
	有 \(f_s'(t)>1/2\)，且 \(f_s(1/2)=1/2\)，所以 \(f_s\) 将 \([0,1)\) 映入自身，并在 \([1/2,1)\) 上恒等。令 \(H_s(c(x,t))=c(x,f_s(t))\)，在领圈外令 \(H_s\) 为恒等。两种表达在领圈参数接近1的区域一致，故拼成光滑同伦，且 \(H_0=\operatorname{id}_M\)。

	由于 \(f_1(0)=\varepsilon>0\)，\(H_1\) 的像完全落在内部，记所得映射为 \(r:M\to\operatorname{Int}M\)。于是 \(j\circ r\simeq\operatorname{id}_M\)。若起点在内部，则领圈参数 \(t>0\)，整个过程中 \(f_s(t)\ge t>0\)，故限制同伦也留在内部，给出 \(r\circ j\simeq\operatorname{id}_{\operatorname{Int}M}\)。领圈已把非紧边界统一参数化为 \([0,1)\)，构造不要求边界紧。
\end{solution}

\begin{exercise}\label{g2:ex:isotopy-complement-normal}
	设 \(K\) 为紧无边界光滑流形，\(N\) 为无边界光滑流形，\(i_t:K\to N\) 为光滑同位。证明 \(N\setminus i_0(K)\) 与 \(N\setminus i_1(K)\) 微分同胚；再证明经 \(i_0,i_1\) 拉回到 \(K\) 的两个法丛同构。说明同位延拓如何同时给出这两个结论。
\end{exercise}
\begin{solution}
	由\cref{g2:thm:isotopy-extension}，存在环境同位 \(\Phi_t\)，使 \(\Phi_t i_0=i_t\)。\(\Phi_1\) 是 \(N\) 的微分同胚，并把 \(i_0(K)\) 恰好送到 \(i_1(K)\)，因此限制到两个补集仍为微分同胚，逆是 \(\Phi_1^{-1}\) 的限制。

	对 \(\Phi_1 i_0=i_1\) 求微分，得
	\[
		\mathrm{d}\Phi_1\circ\mathrm{d}i_0=\mathrm{d}i_1.
	\]
	所以在每个 \(x\in K\) 上，\(\mathrm{d}\Phi_1\) 把 \(\mathrm{d}i_0(\mathrm{T}_xK)\) 同构地映到 \(\mathrm{d}i_1(\mathrm{T}_xK)\)，并在商空间上诱导同构
	\[
		\frac{\mathrm{T}_{i_0(x)}N}{\mathrm{d}i_0(\mathrm{T}_xK)}\longrightarrow
		\frac{\mathrm{T}_{i_1(x)}N}{\mathrm{d}i_1(\mathrm{T}_xK)},\qquad
		[v]\longmapsto[\mathrm{d}\Phi_1(v)].
	\]
	该映射及其逆光滑依赖 \(x\)，故是法丛同构。这里用法商表达最直接；若以度量正交补代表法向量，则须在应用 \(\mathrm{d}\Phi_1\) 后再投影到新的正交补，因为一般微分同胚并不保持所选度量。
\end{solution}

\begin{exercise}\label{g2:ex:moving-compact-support-escape}
	取 \(\chi\in C_c^\infty(\mathbb R)\)，使 \(\chi=1\) 于零的某邻域，\(\operatorname{supp}\chi\subset[-1,1]\)。在实线上定义时间依赖向量场
	\[
		X_t(x)=\begin{cases}
			(1-t)^{-2}\chi\bigl(x-(1-t)^{-1}\bigr)\,\partial_x,&0\le t<1,\\
			0,&t=1.
		\end{cases}
	\]
	证明 \((t,x)\mapsto X_t(x)\) 在 \([0,1]\times\mathbb R\) 上光滑，每个固定时刻的向量场都有紧支集，但从 \(x(0)=1\) 出发的解不能延续到 \(t=1\)。它说明\cref{g2:lem:uniformly-supported-evolution}中的哪个条件不能只逐时刻要求？
\end{exercise}
\begin{solution}
	在 \(t<1\) 处表达显然光滑。给定有限点 \(x_0\)，当 \(t\) 充分接近1且 \(x\) 位于 \(x_0\) 的一个固定有界邻域时，\((1-t)^{-1}\) 与该邻域的距离大于1，故 \(\chi\) 恒为零。于是向量场在 \((1,x_0)\) 的某邻域内恒为零，包括所有阶导数，证明联合光滑性。

	对每个 \(t<1\)，支集包含在紧区间 \([(1-t)^{-1}-1,(1-t)^{-1}+1]\) 内；\(t=1\) 时支集为空。但曲线
	\[
		x(t)=(1-t)^{-1},\qquad0\le t<1,
	\]
	满足 \(x'(t)=(1-t)^{-2}\) 且 \(\chi(x(t)-(1-t)^{-1})=1\)，因此正是所给初值问题的解。它在 \(t\to1\) 时趋于无穷，无法作为实值连续曲线延续。所有时刻的支集并没有包含在一个共同紧集中；逐时刻紧支集允许有效区域逃向无穷，因而不足以保证非自治方程在整个时间区间存在。
\end{solution}

\begin{exercise}\label{g2:ex:compact-fibers-not-proper}
	令 \(E=(S^1\times\mathbb R)\sqcup(S^1\times(0,\infty))\)，令 \(p:E\to\mathbb R\) 在每个分支上都等于第二坐标投影。证明 \(p\) 是满淹没，每条纤维都紧，但 \(p\) 不固有，并且在 \(0\) 附近不局部平凡。
\end{exercise}
\begin{solution}
	两个分支都是无边界二维流形，投影的微分都满射；第一个分支已保证满射性。若 \(t\leq0\)，纤维是一份 \(S^1\)；若 \(t>0\)，纤维是两份 \(S^1\) 的不交并，所以每条纤维均紧。固定 \(z_0\in S^1\)，在第二个分支中取点列 \((z_0,1/j)\)。它位于 \(p^{-1}([0,1])\) 内，却在 \(E\) 中没有收敛子列：第二分支不包含参数零的点，而两个分支都是开闭集，点列也不能收敛到第一分支。因此 \(p^{-1}([0,1])\) 不紧，\(p\) 不固有。\(0\) 的任意邻域都含正参数，纤维的连通分支数由一变二，故不存在局部平凡化。这说明逐条纤维紧不能代替紧集原像紧。
\end{solution}

\begin{exercise}\label{g2:ex:circle-proper-submersion}
	设 \(k\geq2\)，考虑 \(p:S^1\to S^1\)、\(p(z)=z^k\)。证明它是固有满淹没，写出小圆弧上以 \(k\) 个点为纤维的局部平凡化，并证明它没有全局平凡化。
\end{exercise}
\begin{solution}
	源空间紧，故 \(p\) 固有；每个单位复数都有 \(k\) 个互异的单位根，故满射。用角坐标表示为 \(\theta\mapsto k\theta\)，其导数非零，因此是淹没。设圆弧 \(U\) 上选定光滑幅角 \(\alpha:U\to\mathbb R\)，则
	\[
		\{0,1,\ldots,k-1\}\times U\longrightarrow p^{-1}(U),
		\qquad (j,w)\longmapsto\exp\!\left(\frac{\mathrm{i}(\alpha(w)+2\pi j)}{k}\right)
	\]
	是局部平凡化：其像覆盖恰好 \(k\) 个互不相交的开圆弧，在每一份上与 \(p\) 互逆。若存在全局平凡化，则 \(S^1\) 同胚于 \(k\) 份 \(S^1\) 的不交并；前者连通，后者有 \(k\) 个连通分支，矛盾。
\end{solution}

\begin{exercise}\label{g2:ex:proper-submersion-real-line}
	设 \(p:E\to\mathbb R\) 是无边界非空光滑流形上的固有满淹没。证明 \(p\) 有全局平凡化 \(E\cong p^{-1}(0)\times\mathbb R\)。再证明非空紧无边界流形上不存在取值于 \(\mathbb R\) 的处处淹没光滑函数。
\end{exercise}
\begin{solution}
	选取 \(\mathrm{d}p\) 的光滑右逆，令 \(X\) 为常向量 \(1\) 的提升，则 \(\mathrm{d}p(X)=1\)。若 \(e(t)\) 是其积分曲线，则 \(p(e(t))=p(e(0))+t\)。任意有限时间区间中的基底轨道都包含在某个紧区间 \([a,b]\) 内，因此提升留在紧集 \(p^{-1}([a,b])\) 中。逃逸判据说明 \(X\) 完备；这里每个有限时间区间可以使用不同的紧集，不需要 \(E\) 本身紧。

	记其全局流为 \(\varphi^t\)，令 \(F=p^{-1}(0)\)。则
	\[
		\Phi:F\times\mathbb R\longrightarrow E,\quad \Phi(z,t)=\varphi^t(z),
		\qquad \Psi(e)=(\varphi^{-p(e)}(e),p(e))
	\]
	光滑且互逆，并满足 \(p\Phi(z,t)=t\)，这就是所求全局平凡化。最后，紧流形上的实值连续函数取得最大值；在无边界流形的极大点处，光滑函数的微分为零，因而不能处处为淹没。
\end{solution}

\begin{exercise}\label{g2:ex:horizontal-holonomy-product}
	在 \(E=\mathbb R^2/\mathbb Z^2\) 上用 \(([\theta],[\phi])\) 表示坐标，令 \(p:E\to\mathbb R/\mathbb Z\) 为第一坐标投影。给定 \(\alpha\in\mathbb R\)，以 \(X=\partial_\theta+\alpha\partial_\phi\) 张成水平分布。计算沿基底闭路 \(c(t)=[t]\)、\(0\leq t\leq1\) 的水平搬运，并说明这个计算如何区分丛的全局平凡性与所选水平分布的性质。
\end{exercise}
\begin{solution}
	向量场 \(X\) 在整数平移下不变，故在商空间上良定义，且 \(\mathrm{d}p(X)=\partial_\theta\)。经过 \(([0],[\phi_0])\) 的水平提升为
	\[
		e(t)=([t],[\phi_0+\alpha t]).
	\]
	因此绕基底一周后，初始纤维上的搬运映射为 \([\phi]\mapsto[\phi+\alpha]\)。它在 \(\alpha\notin\mathbb Z\) 时不是恒等映射；但 \(p\) 本身就是两个圆的乘积投影，始终全局平凡。故沿闭路搬运的结果还取决于所选水平分布。非平凡的这一搬运映射不能单独证明丛不平凡；这里改取 \(\alpha=0\) 就得到恒等搬运。
\end{solution}
